\chapter{Ophthalmology}

\medterm{Photophobia}-- Abnormal sensitivity or intolerance to light, causing discomfort or the need to avoid bright environments. It may occur with migraine, corneal disease, uveitis, meningitis, retinal disease, medication effects, or neurologic disorders.

\medterm{Ophthalmoscopy}-- Examination of the interior of the eye, especially the retina, optic disc, and retinal blood vessels, using an ophthalmoscope. It helps detect retinal disease, optic-nerve abnormalities, glaucoma, and systemic vascular disease.

\medterm{Watery Eyes}-- Excessive tearing caused by irritation, allergy, infection, dry-eye reflex tearing, or blockage of the tear-drainage system. Persistent unilateral tearing or pain requires ocular assessment.

\medterm{Vision Problems}-- Changes in visual acuity, field, color, or clarity caused by refractive error, eye disease, neurologic disease, or systemic illness. Sudden vision loss is an emergency.

\medterm{Eye Pain}-- Pain in or around the eye caused by surface irritation, infection, trauma, inflammation, glaucoma, or orbital disease. Severe pain with vision loss, photophobia, trauma, or nausea requires urgent assessment.

\medterm{Eye Redness}-- Red discoloration of the eye caused by dilation or inflammation of conjunctival or deeper ocular blood vessels. Causes range from irritation and conjunctivitis to keratitis, uveitis, and acute glaucoma.

\medterm{Eye Floaters}-- Spots, threads, or cobweb-like shapes that move across the visual field, usually caused by opacities in the vitreous. Sudden new floaters, flashes, or a curtain-like shadow require urgent evaluation for retinal tear or detachment. See also Floaters.

\medterm{Nearsightedness}-- A refractive error, also called myopia, in which distant objects appear blurred because light focuses in front of the retina. It is corrected with diverging lenses, contact lenses, or refractive surgery.

\medterm{Farsightedness}-- A refractive error, also called hyperopia, in which light focuses behind the retina when the eye is relaxed. It can cause blurred near vision, eyestrain, and headaches and is corrected with converging lenses or refractive procedures.

\medterm{Interstitial Keratitis}-- Inflammation of the corneal stroma without primary involvement of the corneal surface. It may follow infection, autoimmune disease, or trauma and can cause pain, photophobia, redness, and reduced vision.

\medterm{Optic Glioma}-- A usually slow-growing glioma of the optic nerve or optic pathway, occurring most often in children. It may cause progressive visual loss, optic atrophy, proptosis, or visual-field defects and is associated with neurofibromatosis type 1.

\medterm{Optic Nerve Atrophy}-- Pallor and degeneration of the optic nerve caused by loss of retinal ganglion-cell axons. It produces permanent visual impairment and may result from ischemia, inflammation, compression, glaucoma, trauma, or raised intracranial pressure.

\medterm{Orbital Pseudotumor}-- Noninfectious, idiopathic inflammation of orbital tissues, also called idiopathic orbital inflammation. It can cause painful eye movement, proptosis, eyelid swelling, diplopia, and restricted ocular motility.

\medterm{Occupational Hearing Loss}-- Permanent sensorineural hearing impairment caused by repeated or prolonged exposure to workplace noise or ototoxic hazards. Prevention depends on reducing exposure and using effective hearing protection.

\medterm{Vitreous Hemorrhage}-- Bleeding into the vitreous cavity that causes sudden painless floaters, cobwebs, haze, or loss of vision. Common causes include proliferative diabetic retinopathy, retinal tear or detachment, trauma, and posterior vitreous separation. Dilated examination and ocular ultrasound are used when the retina cannot be visualised; urgent ophthalmologic review is required to exclude retinal detachment.

\medterm{Adie Tonic Pupil}-- Uncommon cause of anisocoria, a benign parasympathetic denervation of the iris sphincter and ciliary muscle, typically unilateral in young women, idiopathicor associated with diabetes, herpes zoster, or autoimmune neuropathy. The affected pupil is enlarged and reacts poorly to light but constricts slowly and tonically to near ("light-near dissociation"). Cholinergic supersensitivity: dilute pilocarpine 0.125\% constricts the Adie pupil but not the normal pupil. Treatment rarely required; reading glasses for accommodation loss.

\medterm{Age-Related Macular Degeneration}-- The leading cause of irreversible central vision loss in individuals over 50 in developed countries, affecting approximately 200 million people worldwide with prevalence increasing with age (10 percent over age 65, 30 percent over age 75). Pathogenesis involves progressive degeneration of the macula with accumulation of drusen (extracellular deposits of lipids, proteins, and complement factors between the retinal pigment epithelium and Bruch membrane) anand loss of retinal pigment epithelium cells, leading to accumulation of lipofuscin and reactive oxygen species that damage photoreceptors and RPE. AMD is classified into early (medium drusen greater than 63 micrometres without pigmentary changes), intermediate (large drusen greater than 125 micrometres with pigmentary abnormalities), and advanced AMD, which is subdivided into dry AMD (geographic atrophy, accounting for 85 to 90 percent of advanced AMD cases) and wet AMD (neovascular or exudative AMD, accounting for 10 to 15 percent of advanced AMD but responsible for 80 to 90 percent of severe vision loss, characterized by the growth of abnormal new blood vessels from the choroid through Bruch membrane into the sub-RPE or subretinal space, which leak fluid, lipid, and blood). Risk factors include increasing age, smoking (the strongest modifiable risk factor, increasing risk 2- to 4-fold), family history, Caucasian race, and genetic susceptibility variants in complement factor H (CFH), ARMS2, and HTRA1 genes.

\medterm{Allergic Conjunctivitis}-- An inflammatory condition of the conjunctiva mediated by IgE-dependent mast cell degranulation in response to environmental allergens, characterized by bilateral ocular itching, tearing, conjunctival injection, and eyelid edema. The condition is classified into two main forms: seasonal allergic conjunctivitis, which occurs during specific pollen seasons and presents with acute symptoms of itching, tearing, and conjunctival injection; and perennial allergic conjunctivitconjunctivitis, which may be less symptomatic but persistent throughout the year, with indoor allergens such as dust mites, mould, and animal dander as common triggers. The ocular surface reaction is characterized by vasodilation of conjunctival blood vessels (causing injection), increased vascular permeability (leading to chemosis), and stimulation of the trigeminal nerve (causing intense itching). Diagnosis is primarily clinical, and treatment consists of topical antihistamine-mast cell stabiliser combination drops (olopatadine, ketotifen, azelastine, epinastine, bepotastine), systemic antihistamines for associated allergic rhinitis, artificial tears, and cold compresses.

\medterm{Alpha-Agonists}-- Alpha-adrenergic agonists are a class of topical medications used in the treatment of glaucoma that lower intraocular pressure through a dual mechanism: reduction of aqueous humor production by the ciliary body through alpha-2 adrenergic receptor stimulation and enhancement of uveoscleral outflow through remodeling of the extracellular matrix. Available agents include brimonidine, apraclonidine, and dipivefrin, with brimonidine being the most commonly used due to its superior tolerability profilprofil. Brimonidine is the most widely prescribed topical alpha-2 agonist for chronic glaucoma therapy, reducing intraocular pressure by 20 to 30 percent. Brimonidine is often used as adjunctive therapy in patients inadequately controlled on prostaglandin analogs and beta-blockers. Common side effects include ocular allergy and conjunctival folliculosis (allergic conjunctivitis develops in up to 10 to 15 percent of patients within 9 to 12 months of initiation), dry mouth, ocular burning and stinging, blurred vision, and fatigue. Apraclonidine is primarily used for acute short-term control of intraocular pressure spikes, such as before, during, and after laser procedures.

\medterm{Amblyopia}-- Amblyopia, commonly known as lazy eye, is a neurodevelopmental disorder of visual acuity that results from abnormal visual experience during the critical period of visual development (birth to approximately eight years of age), leading to reduced best-corrected visual acuity that cannot be attributed to structural ocular pathology. The condition is classified based on the underlying cause: strabismic amblyopia from ocular misalignment, anisometropic amblyopia from unequal refractive errors betwebetween the two eyes, and deprivation amblyopia from obstruction of the visual axis by cataract, ptosis, or corneal opacity. The pathophysiology involves abnormal development of the visual cortex, particularly the lateral geniculate nucleus and primary visual cortex (V1), where the input from the amblyopic eye is suppressed by the cortical dominance of the normal eye. The critical period of susceptibility to amblyopia spans from birth to approximately 8 years of age. The diagnosis is made by detecting a 2-line or greater difference in best-corrected visual acuity between the two eyes. Management consists of correction of any underlying refractive error (the first and most important step), occlusion therapy (patching of the non-amblyopic eye for 2 to 6 hours per day), and pharmacological penalisation (atropine 1 percent drops in the non-amblyopic eye once daily as an alternative to patching).

\medterm{Anisocoria}-- Unequal pupillary diameter exceeding 0.4 mm, physiological in about 20 percent of the population. Pathological causes are localized by which pupil is abnormal and whether asymmetry worsens in bright light (iris sphincter damage, Adie pupil) or dim light (Horner syndrome). Association with ptosis, ophthalmoplegia, headache, or altered mental status demands urgent evaluation for third nerve palsy from intracranial aneurysm.

\medterm{Anterior Chamber Angle}-- The anterior chamber angle is the anatomical region where the cornea and iris meet, and it contains the trabecular meshwork, the primary drainage structure for aqueous humor from the eye. The angle is bounded anteriorly by Schwalbe line (the termination of Descemet membrane), posteriorly by the anterior surface of the iris root, and laterally by the scleral spur and ciliary body band. The trabecular meshwork consists of three regions: the uveal meshwork adjacent to the iris, the corneoscleral memeshwork, which is the outermost layer and the primary site of aqueous humor outflow resistance in primary open-angle glaucoma, and the juxtacanalicular meshwork adjacent to Schlemm canal. The trabecular meshwork is drained by Schlemm canal, a circular, lymphatic-like channel lined by endothelial cells, and from there aqueous humor drains through 25 to 30 collector channels into the episcleral and conjunctival venous system. The anterior chamber angle is evaluated clinically using gonioscopy, which allows direct visualisation of the angle structures and determination of whether the angle is open, narrow, or closed.

\medterm{Anterior Uveitis}-- Anterior uveitis, also known as iritis or iridocyclitis when the ciliary body is involved, is inflammation of the anterior segment of the eye involving the iris and/or ciliary body, and it is the most common form of uveitis. Patients typically present with acute onset of unilateral eye pain, photophobia, blurred vision, and conjunctival injection, and slit-lamp examination reveals cells and flare in the anterior chamber, which represent inflammatory cells and leaked protein from inflamed blood vblood vessels. The etiology of anterior uveitis can be classified as idiopathic (the most common, accounting for 50 to 60 percent of cases), HLA-B27-associated (20 to 30 percent, including patients with ankylosing spondylitis, reactive arthritis, psoriatic arthritis, and inflammatory bowel disease), sarcoidosis (5 to 10 percent, typically with a chronic, bilateral, granulomatous course), and herpes simplex or herpes zoster virus (5 to 10 percent). Treatment consists of topical corticosteroids (prednisolone acetate 1 percent or difluprednate 0.05 percent, with hourly dosing during the acute phase, tapering slowly over weeks to months), cycloplegic agents (homatropine 2 or 5 percent or cyclopentolate 1 percent, to dilate the pupil, relieve ciliary spasm pain, and reduce the risk of posterior synechiae), and systemic therapy for refractory or severe cases.

\medterm{Aqueous Humor}-- A clear, watery fluid that fills the anterior and posterior chambers of the eye, maintaining intraocular pressure and providing nutrients and oxygen to the avascular cornea and lens while removing metabolic waste products. The fluid is produced by the non-pigmented epithelium of the ciliary processes at a rate of approximately 2.5 microliters per minute, through a combination of active secretion (involving carbonic anhydrase and sodium-potassium ATPase), ultrafiltration, and difand diffusion of fluid and electrolytes across the ciliary epithelium, and it drains from the eye through the trabecular meshwork and Schlemm canal into the venous system (conventional outflow pathway, responsible for 80 to 90 percent of outflow resistance and the site of pathology in primary open-angle glaucoma) and through the uveoscleral pathway (unconventional outflow, responsible for 10 to 20 percent of outflow, which is enhanced by prostaglandin analogs).

\medterm{Arcus Senilis (Corneal Arcus)}-- A white, gray, or bluish opaque ring in the peripheral corneal stroma separated from the limbus by a clear zone (lucid interval of Vogt). Caused by lipid (cholesterol) deposition in the peripheral cornea. Common and benign in elderly individuals (>60 years). When present in young patients ($<40$ years, arcus juvenilis), it strongly suggests severe familial hypercholesterolemia or familial combined hyperlipidemia and warrants a lipid panel workup.

\medterm{Argyll Robertson Pupil}-- Small, irregular, unequal pupils that accommodate and converge normally but do not react to light ("prostitute's pupil: accommodates but doesn't react"), caused by neurosyphilis affecting the pretectal midbrain Edinger-Westestphal nucleus pathways, sparing the more ventral pathways for accommodation. Bilateral in tabes dorsalis; may be unilateral with Diabetes mellitus, multiple sclerosis, pineal tumors (Parinaud syndrome), encephalitis, and hereditary ataxias. Workup: syphilis serology (RPR/VDRL with confirmatory FTA-ABS or TPPA) and CSF examination.

\medterm{Astigmatism}-- A refractive error caused by an irregular curvature of the cornea or lens, resulting in light rays being focused at multiple points rather than a single point on the retina, leading to blurred or distorted vision at all distances. The most common form is regular astigmatism, in which the cornea has two principal meridians with different curvatures at right angles to each other, creating a shape similar to a football rather than a basketball. The steepest meridian is called the agaagain or the flattest meridian, and light rays are refracted unequally in the two principal meridians. The diagnosis is made by refraction, keratometry, corneal topography, and autorefraction. Correction is achieved with toric spectacle or contact lenses (which have different powers in the two principal meridians), or refractive surgery including LASIK, PRK, and SMILE, which can reshape the cornea to reduce or eliminate the irregular curvature.

\medterm{Bacterial Conjunctivitis}-- An infection of the conjunctival mucous membrane caused by various bacterial pathogens, presenting with a purulent or mucopurulent discharge that causes the eyelids to become matted together, particularly upon waking, conjunctival injection, and a gritty or foreign body sensation. Acute bacterial conjunctivitis is most commonly caused by Staphylococcus aureus, Streptococcus pneumoniae, and Haemophilus influenzae, and it is usually self-limited but may be treated with topical antibiotics, most commonly a fluoroquinolone (ciprofloxacin, ofloxacin, levofloxacin, moxifloxacin, or gatifloxacin) or polymyxin B-trimethoprim combination, instilled 4 to 6 times daily for 5 to 7 days. Gonorrhoeal conjunctivitis, caused by Neisseria gonorrhoeae, is a sight-threatening emergency presenting with a profuse, purulent discharge, marked chemosis, eyelid edema, and preauricular lymphadenopathy, requiring systemic antibiotic treatment with intramuscular or intravenous ceftriaxone. Chlamydial conjunctivitis, caused by Chlamydia trachomatis serotypes D to K, presents with a chronic, mucopurulent discharge, conjunctival follicles, and keratitis, and requires systemic therapy with azithromycin or doxycycline.

\medterm{Beta-Blockers}-- Beta-adrenergic blocking agents, commonly known as beta-blockers, are topical medications used in the treatment of glaucoma that lower intraocular pressure by reducing the production of aqueous humor by the ciliary body, primarily through blockade of beta-2 adrenergic receptors. Available agents include timolol, levobunolol, metipranolol, and carteolol, typically administered once or twice daily, and they reduce intraocular pressure by approximately twenty to twenty-five percent. Beta-blockers acontraindicated in patients with asthma (due to the risk of bronchospasm), chronic obstructive pulmonary disease, heart block, bradycardia, and congestive heart failure, and they should be used with caution in patients with diabetes mellitus and depression. The alpha-2 agonist brimonidine and the carbonic anhydrase inhibitors dorzolamide and brinzolamide are often used as second-line agents or as adjunctive therapy.

\medterm{Bitemporal Hemianopia}-- Loss of vision in the temporal (lateral) half of both visual fields due to lesions that cross the optic chiasm, damaging the decussating fibers from the nasal retina of each eye (which carry temporal field information). Classic cause: pituitary macroadenoma (typically suprasellar extension) compressing the chiasm from below; craniopharyngioma compressing from above; other suprasellar masses (meningioma, Rathke cleft cyst, aneurysm), traumatic chiasmal injury, demyelination, and glioma. Patients may not spontaneously notice bitemporal loss and present with colliding into doorframes or unrecognized visual field defects; macular split can cause reduced acuity and difficulties with central vision. Diagnosis: confrontation visual field testing, formal perimetry (Humphrey, Goldmann), and MRI pituitary/sellar protocol. Treatment addresses the underlying lesion (transsphenoidal resection for pituitary adenoma, dopamine agonist for prolactinoma, radiation therapy, or observation).

\medterm{Blepharitis}-- Chronic inflammation of the eyelid margins, one of the most common ophthalmic conditions, classified into anterior blepharitis (affecting the eyelash follicle and base) and posterior blepharitis (affecting the meibomian glands). Anterior blepharitis is further subdivided into staphylococcal blepharitis (caused by Staphylococcus aureus colonization, presenting with collarettes and crusting at the base of lashes) and seborrheic blepharitis (associated with seborrheic dermatitis, presenting with grgreasy scales and crusts along the eyelid margin with associated seborrheic dermatitis of the scalp, eyebrows, and nasolabial folds. The pathophysiology involves bacterial colonisation (particularly Staphylococcus epidermidis and Staphylococcus aureus, which produce exotoxins and enzymes that destabilise the tear film), meibomian gland dysfunction (with altered meibum composition, thickened and turbid meibum that obstructs the gland orifices, leading to reduced tear film lipid layer, increased tear evaporation, and tear hyperosmolarity), and ocular surface inflammation. Diagnosis is clinical, based on slit lamp examination of the lid margins, and treatment includes eyelid hygiene (warm compresses for 5 to 10 minutes twice daily followed by lid margin massage and scrubbing), topical antibiotics (bacitracin or erythromycin ointment applied to the lid margins at bedtime, or azithromycin 1 percent solution twice daily), and artificial tears for dry eye symptoms.

\medterm{Bulbar Conjunctiva}-- The bulbar conjunctiva is the portion of the conjunctival membrane that covers the anterior surface of the sclera, extending from the limbus (the junction of the cornea and sclera) to the fornices where it reflects onto the inner surface of the eyelids. It is a thin, transparent, loosely attached membrane through which the underlying white sclera is visible, and it contains blood vessels, lymphatics, and sensory nerve fibers. The bulbar conjunctiva is normally smooth and moist, and its vessels cvessels can become dilated and tortuous in conditions such as conjunctivitis and uveitis, providing important diagnostic information on slit-lamp examination. The bulbar conjunctiva is freely movable over the underlying sclera, to which it is only loosely attached by a thin layer of connective tissue, and this mobility allows for the formation of conjunctival folds or chemosis (edema of the conjunctiva) in inflammatory and allergic conditions. The blood supply to the bulbar conjunctiva arises from two main sources: the anterior ciliary arteries, which are branches off the muscular arteries supplying the extraocular muscles, and the posterior conjunctival arteries from the palpebral arcades.

\medterm{Carbonic Anhydrase Inhibitors}-- A class of medications used in the treatment of glaucoma that lower intraocular pressure by reducing the production of aqueous humor through inhibition of the enzyme carbonic anhydrase, which catalyzes the conversion of carbon dioxide and water to bicarbonate and hydrogen ions, a process essential for the active secretion of bicarbonate into the aqueous humor by the ciliary epithelium. Available topical agents include dorzolamide and brinzolamide, administered tthree times daily. Systemic carbonic anhydrase inhibitors, including acetazolamide and methazolamide, are reserved for acute pressure elevation in angle-closure glaucoma, uveitic glaucoma, and as a temporising measure before surgery, due to the systemic side effects of paresthesias, metallic taste, fatigue, depression, weight loss, and the risk of nephrolithiasis with prolonged use.

\medterm{Cataract}-- A cataract is an opacity or clouding of the crystalline lens that impairs the transmission of light to the retina, resulting in blurred vision, glare, decreased contrast sensitivity, and progressive vision loss, and it is the leading cause of preventable blindness worldwide. The most common form is age-related cataract, which is classified by the location of the opacity within the lens: nuclear sclerosis, characterized by yellowing and hardening of the lens nucleus; cortical cataract, charactericharacterised by radial, spoke-like opacities that extend from the periphery toward the center of the lens; and posterior subcapsular cataract, characterized by granular, plaque-like opacities located at the posterior pole of the lens, just anterior to the posterior capsule, which are particularly visually disabling due to their location in the visual axis and are strongly associated with corticosteroid use, diabetes, trauma, and radiation exposure.

\medterm{Cataracts}-- A clouding of the eye's natural lens that develops gradually with aging, causing blurred vision, glare sensitivity, faded colors, and difficulty seeing at night. They are the leading cause of vision loss worldwide but are treatable with surgical removal of the cloudy lens and replacement with an artificial intraocular lens. Surgery is safe and highly effective, restoring vision in the vast majority of cases.

\medterm{Chalazion}-- A chronic, non-infectious granulomatous inflammation of a meibomian gland caused by duct obstruction. Presents as a firm, painless, well-circumscribed nodule in the eyelid, usually away from the lid margin. Different from a stye (which is acute and tender). Most resolve spontaneously over weeks to months. Treatment: warm compresses, lid massage. For persistent chalazia: intralesional corticosteroid injection or incision and curettage (I\&C).

\medterm{Cherry-Red Spot}-- A striking fundoscopic finding in which the macula appears bright red against a pale, ischemic, or lipid-laden surrounding retina. Mechanism: the fovea centralis lacks inner retinal layers and receives blood supply from the underlying choroid, preserving its natural red color while adjacent edematous retina turns opaque white. Classic causes: Central Retinal Artery Occlusion (CRAO — sudden painless monocular vision loss), Tay-Sachs Disease (hexosaminidase A deficiency), Niemann-Pick Disease (sphingomyelinase deficiency), and Gaucher Disease.

\medterm{Choroidal Neovascularization}-- The pathological growth of new, fragile blood vessels from the choroid through Bruch membrane into the subretinal or sub-retinal pigment epithelium space, and it is the hallmark of wet age-related macular degeneration and the primary cause of severe central vision loss in this condition. The neovascular vessels are stimulated by vascular endothelial growth factor, which is produced in response to hypoxia, inflammation, and oxidative stress in the aging retina, andand they are fragile and prone to leakage and bleeding, leading to sudden central vision loss from hemorrhage as well as gradual vision loss from exudation and fibrosis. The diagnosis of choroidal neovascularization is confirmed by fluorescein angiography (which shows the classic lacy, petaloid hyperfluorescence in the subretinal space with early leakage and late staining in the classic type) and optical coherence tomography (which shows subretinal or sub-RPE fluid, intraretinal cystic spaces, and a hyperreflective lesion above or below the RPE).

\medterm{Color Blindness}-- Inability or decreased ability to perceive color differences under normal lighting conditions. Most commonly X-linked recessive, affecting ~8\% of males and 0.5\% of females. Types: deuteranopia (green deficiency, most common), protanopia (red deficiency), tritanopia (blue-yellow deficiency, rare), and achromatopsia (complete color blindness, very rare). Most congenital color blindness is non-progressive and does not affect visual acuity. Diagnosis: Ishihara color plates. No treatment available; patients adapt with compensatory strategies.

\medterm{Color Vision Deficiency}-- Color vision deficiency, commonly known as color blindness, is a condition in which the ability to perceive and discriminate certain colors is reduced or absent, most commonly affecting the perception of red and green hues, and it is caused by abnormalities in the cone photoreceptors responsible for color vision. The most common form is red-green color deficiency, which is an X-linked recessive condition affecting approximately eight percent of males and 0.5 percent of females, caused by mutatiomutations in the opsin genes (OPN1LW and OPN1MW on the X chromosome encoding the long-wavelength and middle-wavelength cone photopigments, respectively, which cause the common red-green color deficiencies through deletions, fusions, or point mutations). Blue-yellow color deficiency (tritanopia) is much rarer and is caused by mutations in the OPN1SW gene on chromosome 7 encoding the short-wavelength cone photopigment, and it is inherited as an autosomal dominant condition. Complete color blindness (achromatopsia) is a rare autosomal recessive condition characterized by the absence of all three types of cone function, resulting in severe photophobia, nystagmus, and poor visual acuity.

\medterm{Congenital Cataract}-- Opacification of the crystalline lens present at birth, a leading cause of treatable childhood blindness worldwide. Causes: genetic (autosomal dominant most common, mutations in crystallin genes or PAX6), infectious (congenital rubella syndrome, toxoplasmosis, CMV, syphilis, HSV), metabolic (galactosemia, hypoglycemia, Wilson disease), and syndromic (Down syndrome, Lowe syndrome, Marfan syndrome). Bilateral cataracts are more commonly genetic or metabolic; unilateral cataracts suggest structural eye anomalies or asymmetric intrauterine infection. Presents with leukocoria (white pupillary reflex), nystagmus, strabismus, or poor visual tracking. Screening: red reflex examination (absent or abnormal). Diagnosis: slit-lamp examination after dilation. Treatment: early surgical extraction (within 4-6 weeks for dense bilateral cataracts to prevent deprivation amblyopia), primary or secondary intraocular lens implantation (depending on age), and immediate postoperative aphakic correction with contact lenses or glasses. Lifelong monitoring for glaucoma, strabismus, and amblyopia.

\medterm{Congenital Glaucoma}-- A rare form of glaucoma (1 in 10,000-20,000 births) caused by developmental anomalies of the trabecular meshwork and anterior chamber angle, impairing aqueous humor outflow. Most cases are sporadic; familial forms are autosomal recessive (CYP1B1 gene mutations). Presents in the first year of life with the classic triad: epiphora (excessive tearing, often the earliest sign), photophobia, and blepharospasm. Corneal edema causes corneal enlargement (buphthalmos, or ox eye, from elevated IOP stretching the immature collagen), Haab striae (linear breaks in Descemet membrane), and corneal clouding. Increased axial length leads to progressive myopia. Diagnosis: examination under anesthesia — elevated intraocular pressure, increased corneal diameter (>12 mm by 1 year), and optic nerve cupping (reversible in early stages). Treatment: surgical angle surgery (goniotomy or trabeculotomy) is first-line (success rate 80-90\% if performed early). Trabeculectomy or drainage implants for failed angle surgery. Lifelong IOP monitoring is essential.

\medterm{Conjunctiva}-- The conjunctiva is a thin, transparent, mucous membrane that lines the inner surface of the eyelids (palpebral conjunctiva) and covers the anterior surface of the sclera (bulbar conjunctiva), forming a continuous sheet that protects the eye and facilitates smooth eyelid movement. The palpebral conjunctiva is further divided into the tarsal conjunctiva, which is firmly adherent to the tarsal plates of the eyelids, and the forniceal conjunctiva, which forms loose folds at the junction between the tarsal and forniceal conjunctiva. The conjunctiva is composed of a non-keratinising stratified columnar epithelium with goblet cells that secrete mucin, a critical component of the tear film that stabilises the tear film by lowering its surface tension, and an underlying substantia propria, a layer of loose connective tissue containing blood vessels, lymphatics, nerves, and immune cells. The conjunctiva is a highly vascularised, immunologically active tissue that protects the eye from pathogens and environmental insults.

\medterm{Conjunctivitis}-- Inflammation of the conjunctiva, the thin, transparent mucous membrane lining the inner surface of the eyelids and the anterior surface of the sclera. Conjunctivitis is the most common cause of red eye and is classified by etiology. Infectious conjunctivitis: viral (most common overall, caused by adenovirus serotypes 3, 4, 7, 8, 19, 37, presenting with watery discharge, preauricular lymphadenopathy, and highly contagious, often involving both eyes with one eye affected first), bacterial (purulenpurulent discharge, often unilateral, treated with topical antibiotics), and chlamydial (inclusion conjunctivitis, with a chronic course, mucopurulent discharge, and follicles, requiring systemic antibiotics). Infectious conjunctivitis is distinguished from non-infectious causes: allergic conjunctivitis (intense itching, bilateral, seasonal or perennial, with diffuse conjunctival injection and chemosis, treated with antihistamine-mast cell stabiliser drops) and chemical or irritant conjunctivitis (acute exposure to chemicals or irritants, with immediate onset of redness, tearing, and pain, treated with copious irrigation and lubrication).

\medterm{Corneal Abrasion}-- Loss of corneal epithelium caused by mechanical trauma — contact with fingernail, makeup brush, foreign body, contact lens overwear, dry eye, sports injury, paper edge, or surgical instrumentation. Most common eye injury presenting to emergency departments. Clinical: sudden-onset sharp ocular pain, foreign body sensation, photophobia, tearing, blepharospasm, blurred vision; pain typically worse with blinking and eye movement (eyelid rubbing exposed defect), and may awaken patient from sleep (cornea dries overnight). Diagnosis: focal fluorescein uptake (defect stains bright green under cobalt blue; vertical abrasions suggest foreign body trapped under upper lid — evert and inspect), reduced corneal sensation in area; slit-lamp examination to assess depth, infiltrate, and underlying tear film/dry eye; document pupil reaction (an AbP carry implant). Treatment: topical antibiotic prophylaxis (chloramphenicol, erythromycin ointment, fluoroquinolone drops for contact lens wear — including Pseudomonas coverage); cycloplegia (cyclopentolate 1\% or homatropine 2\%) for pain from ciliary spasm; oral NSAIDs/acetaminophen for comfort; \emph{avoid} pressure patching (does not accelerate healing, increases Pseudomonas risk in contact-lens wearers, hides complications); contact lens discontinuation until healed; follow-up in 24-48 h to confirm healing. Differential — corneal ulcer (infiltrate, hypopyon), keratitis, foreign body, recurrent erosion (history of abrasion, recurrent early-morning pain).

\medterm{Corneal Cross-Linking}-- Corneal collagen cross-linking is a minimally invasive therapeutic procedure that strengthens the corneal stroma by creating additional chemical bonds between collagen fibrils, thereby stiffening the cornea and halting or slowing the progression of corneal ectatic disorders, particularly keratoconus and post-refractive surgery ectasia. The procedure involves debriding the central corneal epithelium, applying riboflavin (vitamin B2) drops to the cornea for approximately thirty minutes to saturatesaturate the corneal stroma with riboflavin, and then exposing the cornea to ultraviolet A light (365 nm, 3 mW/cm2) for 30 minutes, which causes the riboflavin to undergo a photochemical reaction that generates reactive oxygen species that induce new covalent cross-links between collagen molecules, increasing the biomechanical rigidity of the cornea by 3 to 4 times. The epithelium-on (epi-on) and epithelium-off (epi-off) techniques exist, with epi-off being the traditional Dresden protocol and still considered the gold standard.

\medterm{Corneal Ulcer}-- A corneal ulcer is a defect in the corneal epithelium with associated stromal inflammation and necrosis, typically caused by infectious or non-infectious mechanisms, and it is a sight-threatening condition that requires prompt diagnosis and treatment. Bacterial corneal ulcers, the most common infectious type, are associated with contact lens wear, ocular surface disease, and corneal trauma, and common causative organisms include Pseudomonas aeruginosa, Staphylococcus aureus, and Streptococcus pnStreptococcus pneumoniae. Fungal corneal ulcers, more common in tropical regions and associated with agricultural trauma, contact lens wear, and topical corticosteroid use, are caused by filamentous fungi (Fusarium, Aspergillus) and yeasts (Candida), and they present with a dry, elevated infiltrate with a feathery border and satellite lesions. Viral corneal ulcers, most commonly caused by herpes simplex virus and varicella-zoster virus, present with a characteristic dendritic epithelial defect on fluorescein staining, and they are treated with topical or oral antivirals (acyclovir, ganciclovir, valacyclovir) and cycloplegia.

\medterm{Cortical Cataract}-- A type of age-related lens opacity that begins in the peripheral lens cortex as wedge-shaped or spoke-like opacities that extend centripetally toward the visual axis, caused by disruption of the normally orderly arrangement of lens fibers and the accumulation of fluid between them. The opacities appear as white, wedge-shaped streaks in the lens cortex on slit-lamp examination, with the narrow end of the wedge pointing toward the lens center, and they may progress centrally ttoward the visual axis, eventually involving the central visual axis and causing significant visual impairment. The cortical spokes may form a wedge-shaped opacity with a clear center (cuneiform cataract) or may extend centrally to create a radial, star-shaped pattern. Patients with cortical cataract often complain of glare and difficulty with night vision.

\medterm{Cup-to-Disc Ratio}-- The cup-to-disc ratio is a clinical measurement used to assess the severity of glaucomatous optic nerve damage, comparing the diameter of the optic cup (the central depression within the optic nerve head where neural tissue is absent) to the diameter of the optic disc (the entire optic nerve head). A normal cup-to-disc ratio ranges from 0.1 to 0.3, while progressive enlargement of the cup, known as cupping, indicates loss of retinal ganglion cell axons and is a hallmark of glaucomatous optic neuneuropathy. The cup-to-disc ratio is assessed during fundoscopic examination of the optic nerve head, and its measurement is essential for the diagnosis and monitoring of glaucoma, with a focal or diffuse enlargement of the cup (increased cup-to-disc ratio), asymmetry of cup-to-disc ratio between the two eyes (a difference of greater than 0.2 is suspicious for glaucoma), and notching or thinning of the neuroretinal rim being characteristic of glaucomatous optic nerve damage.

\medterm{Dacryoadenitis}-- Inflammation of the lacrimal gland, presenting as tender swelling in the superolateral aspect of the upper eyelid with characteristic S-shaped ptosis. It can be acute or chronic in presentation. Acute dacryoadenitis is most commonly caused by viral infection (Epstein-Barr virus, mumps, adenovirus, herpes zoster, enterovirus), with bacterial causes (Staphylococcus aureus, Streptococcus, Haemophilus influenzae, gonococcus) less common but more severe. Chronic dacryoadenitis is associated with autoautoimmune conditions including sarcoidosis, rheumatoid arthritis, Sjogren syndrome, granulomatosis with polyangiitis, systemic lupus erythematosus, IgG4-related disease, and thyroid eye disease. Diagnosis is clinical, with slit-lamp examination showing an enlarged, erythematous lacrimal gland. Imaging with CT or MRI may be needed to confirm the diagnosis and rule out orbital mass lesions. Treatment of acute viral dacryoadenitis is supportive, while bacterial dacryoadenitis requires systemic antibiotics. Chronic dacryoadenitis requires treatment of the underlying systemic condition.

\medterm{Dacryocystitis}-- Inflammation of the lacrimal sac, most commonly caused by obstruction of the nasolacrimal duct leading to stagnation of tears and secondary bacterial infection (Staphylococcus aureus, Streptococcus pneumoniae, Haemophilus influenzae). Presents with pain, swelling, erythema, and tenderness over the medial canthus, with epiphora (excessive tearing) and purulent discharge expressed from the punctum. Acute dacryocystitis requires systemic antibiotics and warm compresses; chronic cases may require dacryocystorhinostomy (DCR) to bypass the obstruction.

\medterm{Dendritic Ulcer}-- A characteristic branching (dendritic) epithelial ulcer of the cornea caused by herpes simplex virus (HSV) infection. It is the most common infectious cause of corneal blindness in developed countries. Presents with unilateral red eye, photophobia, tearing, foreign body sensation, and decreased vision. Slit-lamp examination with fluorescein staining shows a branching linear ulcer with terminal bulbs (dendrite). Treatment: topical antivirals (trifluridine drops, ganciclovir gel) or oral antivirals (acyclovir, valacyclovir). Topical corticosteroids are contraindicated in active epithelial disease.

\medterm{Detached Retina}-- Separation of the neurosensory retina from the underlying retinal pigment epithelium (RPE), a medical emergency requiring prompt surgical repair to prevent permanent vision loss. Types: rhegmatogenous (most common--a tear or hole in the retina allows vitreous fluid to accumulate in the subretinal space, associated with posterior vitreous detachment, high myopia, trauma, and aphakia), tractional (vitreoretinal adhesions from proliferative diabetic retinopathy, retinopathy of prematurity, or penetrating trauma pull the retina away), and exudative (fluid accumulation beneath the retina without a tear, from inflammation, tumors, or choroidal disorders). Presents with sudden onset of flashes (photopsia), floaters, and a curtain-like or shadowy visual field defect. Treatment: pneumatic retinopexy, scleral buckle, or pars plana vitrectomy with gas or silicone oil tamponade.

\medterm{Diabetic Macular Edema}-- The accumulation of fluid within the macula due to breakdown of the blood-retinal barrier in diabetic retinopathy, resulting in thickening of the macula and central vision loss, and it can occur at any stage of diabetic retinopathy, from mild non-proliferative to proliferative disease. The pathogenesis involves hyperglycemia-induced damage to retinal vascular endothelial cells and pericytes, leading to increased vascular permeability, leakage of plasma components into tinto the retinal tissue, and thickening of the macula. The severity of DME is graded on clinical examination based on the distance of the retinal thickening from the fovea. The diagnosis of DME is confirmed and quantified by optical coherence tomography. The treatment of center-involving DME is primarily with intravitreal anti-VEGF agents (ranibizumab, aflibercept, bevacizumab, and faricimab), which have revolutionised the management of DME.

\medterm{Diplopia (Double Vision)}-- The perception of two images of a single object, classified by onset, duration, and whether it resolves with covering one eye. Monocular diplopia (persists with one eye covered): caused by ocular surface abnormalities (dry eye, irregular astigmatism, cataract, corneal scar, lens dislocation)--usually benign. Binocular diplopia (resolves with either eye covered): results from misalignment of the visual axes due to cranial nerve palsy (CN III, IV, VI), neuromuscular junction disease (myasthenia gravis), restrictive orbital disease (thyroid eye disease), or strabismus. Evaluation: cover test in nine gaze positions, forced duction testing, and neuroimaging if nerve palsy suspected.

\medterm{Drusen}-- Yellowish-white deposits of extracellular material that accumulate beneath the retinal pigment epithelium between the retinal pigment epithelium basement membrane and Bruch membrane, and they are the hallmark clinical finding in dry age-related macular degeneration. Drusen are composed of lipids, complement components, inflammatory mediators, and other proteins, and they represent localized abnormalities of the retinal pigment epithelium basement membrane. Small hard drusen (less thanthan 63 micrometres in diameter) are common in the aging retina and are not associated with an increased risk of progression to advanced AMD, while medium (63 to 125 micrometres) and large (greater than 125 micrometres) drusen, particularly when associated with pigmentary abnormalities, are the hallmark of early and intermediate AMD and are associated with a significantly increased risk of progression to advanced AMD.

\medterm{Dry Eye Disease}-- Dry eye disease, also known as keratoconjunctivitis sicca, is a multifactorial disorder of the ocular surface characterized by loss of homeostasis of the tear film, resulting in symptoms of ocular discomfort, visual disturbance, and tear film instability, with potential damage to the ocular surface. The disease is classified into two major subtypes: aqueous-deficient dry eye, caused by decreased lacrimal gland production of the aqueous layer of the tear film, as seen in Sjogren syndrome and aginaging in non-Sjogren dry eye), and evaporative dry eye, caused by meibomian gland dysfunction (the most common cause of evaporative dry eye, accounting for 80 percent of dry eye cases in some series), where a reduction in the quantity or quality of the lipid layer of the tear film leads to increased tear evaporation and tear hyperosmolarity. The diagnosis of dry eye disease is made by a combination of patient symptom assessment using validated questionnaires and objective clinical tests including the Schirmer test, tear film breakup time, corneal and conjunctival staining, and evaluation of the meibomian glands.

\medterm{Dry Eye Syndrome}-- A multifactorial disease of the ocular surface characterized by tear film instability, hyperosmolarity, and inflammation. Aqueous-deficient dry eye: Sjogren syndrome (autoimmune), non-Sjogren (lacrimal gland dysfunction, age). Evaporative dry eye: meibomian gland dysfunction (most common), incomplete blink, environmental factors. Symptoms: dryness, grittiness, burning, photophobia, blurred vision that improves with blinking. Diagnosis: Schirmer test, tear breakup time, corneal fluorescein staining, meibomian gland evaluation. Treatment: artificial tears (preservative-free), warm compresses, lid hygiene, topical cyclosporine, punctal plugs, and omega-3 supplements.

\medterm{Endophthalmitis}-- A severe, sight-threatening intraocular inflammation involving the vitreous cavity, aqueous humor, and adjacent structures, representing an ophthalmologic emergency. It is classified as exogenous (most common, caused by direct inoculation from intraocular surgery particularly cataract extraction, intravitreal injections, penetrating trauma, or corneal infection) or endogenous (metastatic seeding from a distant infection source, more common in immunocompromised patients, endocarditis, intravenousintravenous drug use, indwelling catheters, and septicemia, with the most common causative organisms being Streptococcus species, Staphylococcus aureus, and gram-negative bacilli). The condition presents with rapid onset of severe, progressive eye pain, decreased visual acuity, conjunctival injection, chemosis, corneal edema, hypopyon (the hallmark of endophthalmitis, present in 80 to 90 percent of cases), and vitreous inflammation. The Endophthalmitis Vitrectomy Study established the management protocol: immediate intravitreal injection of antibiotics (vancomycin 1 mg/0.1 mL for gram-positive coverage and ceftazidime 2.25 mg/0.1 mL for gram-negative coverage), with vitreous tap for culture, and vitrectomy for patients with light perception vision only.

\medterm{Entropion}-- An eyelid malposition in which the eyelid margin rolls inward, causing the eyelashes and skin to rub against the ocular surface, resulting in corneal irritation, conjunctival injection, tearing, photophobia, and a foreign body sensation. The condition most commonly affects the lower eyelid and is classified into four types. Involutional entropion (most common, age-related, due to horizontal laxity from canthal tendon weakening, dehiscence or attenuation of the lower eyelid retractors, and overrioverriding of the preseptal orbicularis oculi muscle, allowing the lower eyelid to roll inward), cicatricial entropion (caused by scarring and contraction of the posterior lamella, most commonly following chemical burns, Stevens-Johnson syndrome, ocular cicatricial pemphigoid, or chronic trachoma), congenital entropion (rare, due to hypertrophic marginal bundle of the orbicularis oculi muscle), and spastic entropion (temporary, due to sustained orbicularis oculi contraction). The diagnosis is clinical, based on slit-lamp examination.

\medterm{Epiphora}-- Excessive tearing caused by impaired tear drainage or reflex hypersecretion. Causes: nasolacrimal duct obstruction (most common in infants and elderly), punctal stenosis, eyelid malposition (ectropion, entropion), and reflex tearing from dry eye or ocular surface irritation. Congenital NLDO: presents with tearing and mucoid discharge in infants; most resolve spontaneously by 12 months. Diagnosis: dye disappearance test, lacrimal irrigation, dacryocystography. Treatment: lacrimal massage for infants; probing and irrigation; dacryocystorhinostomy (DCR) for adults with structural obstruction.

\medterm{Episcleritis}-- A benign, self-limited inflammatory condition of the episclera, the thin layer of connective tissue between the sclera and conjunctiva. It is most common in young to middle-aged women and is typically idiopathic, though it is associated with systemic inflammatory conditions in approximately 30 percent of cases (rheumatoid arthritis, systemic lupus erythematosus, Sjogren syndrome, granulomatosis with polyangiitis, inflammatory bowel disease, psoriatic arthritis, and ankylosing spondylitis). The pThe patient presents with acute or subacute onset of unilateral or bilateral eye redness, a dull ache or irritation, and sectoral injection of the episcleral vessels (which appear as bright red, radially oriented, tortuous engorged vessels that blanch with topical phenylephrine 2.5 percent, distinguishing episcleritis from scleritis). The diagnosis is clinical. The condition is self-limited and resolves spontaneously within 2 to 3 weeks in most cases.

\medterm{Esotropia}-- A type of strabismus in which one eye deviates inward (toward the nose). Congenital/infantile esotropia: large-angle, constant, onset before 6 months; requires early surgical correction. Accommodative esotropia: associated with hyperopia, onset age 2-3 years; treatment with corrective glasses improves alignment. Intermittent esotropia can be controlled with fusion. Complications: amblyopia (lazy eye) and loss of binocular vision. Diagnosis: cover-uncover test, alternate cover test, cycloplegic refraction. Management: glasses, patching, prism therapy, and strabismus surgery.

\medterm{Exophthalmos}-- Abnormal protrusion of the eyeball (globe) from the orbit, most commonly caused by thyroid eye disease (Graves ophthalmopathy--an autoimmune condition associated with hyperthyroidism). The proptosis results from inflammation, edema, and fibrosis of the extraocular muscles and proliferation of orbital fat within the fixed bony orbit. Bilateral exophthalmos is classic for Graves disease. Unilateral exophthalmos suggests orbital tumor (hemangioma, meningioma, optic nerve glioma, lymphoma), orbital pseudotumor, cellulitis, or carotid-cavernous fistula. Measurement: Hertel exophthalmometry (normal <20 mm, or asymmetry >2 mm). Symptoms: proptosis, lid retraction, lid lag (lagophthalmos), exposure keratopathy, diplopia, and optic nerve compression. Treatment: glucocorticoids, orbital decompression surgery, and management of the underlying thyroid disease.

\medterm{Exotropia}-- A type of strabismus in which one eye deviates outward (toward the temple). Intermittent exotropia: most common type, often noticed when tired or daydreaming, can be controlled with fusion. Constant exotropia: requires treatment to prevent amblyopia. Congenital exotropia: rare, requires early evaluation. Diagnosis: cover test, angle measurement. Treatment: part-time patching, orthoptic exercises (convergence exercises), prism glasses, and surgical correction for persistent deviations.

\medterm{Eye Disorders}-- The sensory ability to perceive light and interpret visual information processed by the eyes and brain, enabling recognition of shapes, colors, depth, and motion. See \hyperlink{Vision}{Vision}.

\medterm{Eye Health}-- The maintenance of optimal vision and prevention of ocular disease through regular eye examinations, protective measures, and timely treatment. Common conditions affecting eye health include refractive errors, cataracts, glaucoma, diabetic retinopathy, and age-related macular degeneration. Protective strategies include wearing UV-blocking sunglasses, managing systemic health conditions, and avoiding eye strain.

\medterm{Flashes}-- Perception of flickering light or lightning-like streaks from mechanoelectrical stimulation of the retina rather than an external source. The leading cause is posterior vitreous detachment, often with floaters; retinal tears and detachment, migraine aura, and optic neuritis also produce flashes. New flashes with floaters or a curtain-like field defect require urgent examination to exclude retinal break or detachment.

\medterm{Floaters}-- Mobile spots, threads, cobwebs, or rings in the visual field caused by vitreous opacities casting shadows on the retina. Most arise from posterior vitreous detachment, an age-related separation of the vitreous gel from the retina; hemorrhage, vitritis, and asteroid hyalosis are other causes. Sudden new floaters with flashes or field loss require prompt retinal evaluation to exclude tear or detachment.

\medterm{Fluorescein Angiography}-- An ophthalmological diagnostic procedure that involves intravenous injection of sodium fluorescein dye and sequential photography of the retinal vasculature using a specialized fundus camera fitted with excitation and barrier filters, providing detailed information about retinal blood flow, vascular permeability, and retinal pigment epithelium integrity. The dye is excited by blue light at approximately 490 nanometers and emits yellow-green fluorescence at approximatelapproximately 520 to 530 nanometres. The procedure involves a series of timed images captured in rapid succession following the injection: the early arterial phase (10 to 15 seconds), the arteriovenous phase (15 to 20 seconds), the venous phase (20 to 30 seconds), and the late phase (5 to 10 minutes). The normal fluorescein angiogram shows the retinal vessels fluoresce brightly against the dark background of the choroidal fluorescence.

\medterm{Fornix}-- The conjunctival fornix is the redundant, dome-shaped fold of conjunctiva where the palpebral conjunctiva transitions to the bulbar conjunctiva, forming a deep recess superiorly (superior fornix), inferiorly (inferior fornix), and medially and laterally (medial and lateral fornices). The fornices are important anatomical landmarks for contact lens fitting, conjunctival sampling for diagnosis of chlamydial and viral infections, and as reservoirs for topical medications applied to the eye. The supThe superior fornix is the deepest, extending approximately 8 to 10 mm above the superior corneal limbus, while the inferior fornix is approximately 6 to 8 mm deep. The fornices are also the site of lymphatic drainage of the conjunctiva, which drains to the preauricular and submandibular lymph nodes, and they contain accessory lacrimal glands (the glands of Krause and Wolfring) that contribute to the basal secretion of the aqueous component of the tear film.

\medterm{Fovea}-- A tiny depression in the center of the retina that provides the sharpest and most detailed visual acuity.

\medterm{Fovea Centralis}-- Central macular pit (approximately 1.5 mm) with the highest cone density (approximately 200,000 per square millimeter) and no rods, providing the sharpest visual acuity and color vision. See \hyperlink{Anatomy}{anatomy}.

\medterm{Fundus Photography}-- The documentation of the appearance of the retina, optic nerve, and retinal vasculature using a specialized camera that illuminates and photographs the posterior segment of the eye, providing a permanent, objective record of retinal findings that can be used for comparison over time, patient education, and telemedicine consultations. Standard fundus photography captures images of the optic disc, macula, and major retinal vessels, and widefield fundus photography can image uup to 200 degrees of the retina in a single image. Fundus photography can be performed with or without pupillary dilation. The images are taken in color, red-free (using green light to enhance the visibility of retinal vessels and haemorrhages), and autofluorescence (using blue light to excite lipofuscin in the RPE).

\medterm{Geographic Atrophy}-- The advanced form of dry age-related macular degeneration, characterized by well-demarcated, progressively enlarging areas of complete loss of the retinal pigment epithelium, choriocapillaris, and overlying photoreceptors in the macula, resulting in irreversible central vision loss. The atrophic areas appear as well-defined, depigmented regions on fundus examination, with visible underlying choroidal vessels, and they typically begin in the parafoveal region and expand centcentrally over time, leading to a progressive, bilateral, irreversible, and painless loss of central vision, with the preservation of peripheral vision until late in the disease process. The diagnosis of geographic atrophy is based on fundoscopic examination (showing well-demarcated, depigmented areas with visible choroidal vessels), fundus autofluorescence (showing areas of absent autofluorescence), and optical coherence tomography (showing the loss of the RPE and outer retinal layers).

\medterm{Glaucoma}-- A group of progressive optic neuropathies characterized by the degeneration of retinal ganglion cells and their axons, resulting in characteristic optic nerve head changes and visual field loss, with elevated intraocular pressure being the most important modifiable risk factor. Primary open-angle glaucoma, the most common form, involves gradual obstruction of the trabecular meshwork with chronically elevated intraocular pressure, progressive cupping of the optic nerve head, and characharacteristic visual field defects (paracentral scotomas, arcuate scotomas, and nasal steps in the early stages, progressing to central vision loss and tunnel vision in advanced disease). Other forms of glaucoma include normal-tension glaucoma (glaucomatous optic neuropathy and visual field loss with intraocular pressure consistently within the normal range), primary angle-closure glaucoma (caused by pupillary block, leading to an acute, painful, red eye with fixed, mid-dilated pupil, corneal edema, and elevated IOP), and secondary glaucomas resulting from other ocular or systemic disease.


\medterm{Goldmann Tonometry}-- Goldmann applanation tonometry is the gold standard method for measuring intraocular pressure, the most important modifiable risk factor for glaucoma, and it involves applanating (flattening) a fixed area of the cornea using a biprism attached to a slit lamp, measuring the force required to achieve this applanation. The principle is based on the Imbert-Fick law, which states that the pressure inside a sphere is equal to the force needed to flatten a given area of its surface, and the Goldmann tothe Goldmann tonometer is calibrated to measure the intraocular pressure based on the force required to applanate a 3.06 mm diameter area of the central cornea. The Goldmann tonometer is mounted on a slit lamp, topical anesthesia is applied to the cornea, and fluorescein dye is added to the tear film to stain the tear meniscus.

\medterm{Gonioscopy}-- An ophthalmological examination technique that allows direct visualization of the anterior chamber angle, the anatomical region where the cornea and iris meet and where the trabecular meshwork is located, and it is essential for distinguishing between open-angle and angle-closure forms of glaucoma. The procedure uses a specialized goniolens placed on the surface of the eye that overcomes total internal reflection at the corneal-air interface, allowing light to enter and illuminate the angle structures for gonioscopic evaluation. Gonioscopy is performed using either a direct goniolens (such as the Koeppe lens, providing a direct view of the angle) or an indirect goniolens (such as the Goldmann three-mirror lens or the Zeiss four-mirror lens, providing an indirect, mirrored view of the angle using mirrors at various angles).

\medterm{Hemianopia}-- Loss of vision in half the visual field of one or both eyes. Homonymous hemianopia affects the same side in both eyes and localizes to a contralateral lesion posterior to the optic chiasm, most often stroke or tumor. Bitemporal hemianopia affects both temporal fields and localizes to the chiasm, classically from pituitary adenoma or craniopharyngioma.

\medterm{Hordeolum (Stye)}-- A common infection of the eyelid. See \hyperlink{Stye}{Stye}.

\medterm{Hyperopia}-- Hyperopia, commonly known as farsightedness, is a refractive error in which parallel light rays from distant objects are focused posterior to the retina rather than directly on it, resulting in blurred near vision that may be accompanied by blurred distance vision in moderate to high degrees. The condition is caused by an insufficient axial length of the globe relative to the refractive power of the cornea and lens, or by a flat corneal curvature, and it is corrected with convex (plus-powered) cconvex (plus-powered) converging lenses (spectacles or contact lenses) that move the focal point of light from behind the retina back onto the retina. Hyperopia is quantified by the refractive error in dioptres, with the degree classified as low (less than +2.00 D), moderate (+2.00 D to +5.00 D), and high (greater than +5.00 D).

\medterm{Hypertensive Retinopathy}-- Retinal changes caused by chronic or acute hypertension. Grading (Keith-Wagener-Barker): Grade I: generalized arteriolar narrowing; Grade II: arteriovenous nicking; Grade III: cotton-wool spots, flame hemorrhages, hard exudates; Grade IV: papilledema (malignant hypertension). Severe hypertension can cause retinal artery occlusion, vitreous hemorrhage, and vision loss. Management: blood pressure control; ophthalmology referral for advanced grades.

\medterm{Hyphema}-- Bleeding into the anterior chamber of the eye, between the cornea and iris, typically resulting from blunt ocular trauma that tears the blood vessels of the iris or ciliary body. Less common causes include spontaneous hyphema (from iris neovascularization in diabetes, retinal vein occlusion, or ocular ischemic syndrome), intraocular surgery (post-cataract or glaucoma surgery), and coagulopathy (anticoagulant or antiplatelet therapy, hemophilia, von Willebrand disease, thrombocytopenia). The blooThe blood in the anterior chamber layers dependently due to gravity, forming a visible level (the hyphema) that is graded based on the height of the blood level relative to the total anterior chamber height. The immediate management includes strict bed rest with the head elevated at 30 to 45 degrees, shielding the eye with a protective patch, and cycloplegia with atropine 1 percent or homatropine 5 percent.

\medterm{Hypopyon}-- A layering of inflammatory cells (purulent exudate/leukocytes) in the anterior chamber of the eye, forming a visible white-to-yellow fluid level in the lower portion of the cornea. Caused by severe anterior uveitis (HLA-B27 spondyloarthropathies, Behçet disease), infectious keratitis (bacterial/fungal corneal ulcers), or endophthalmitis. Ophthalmic emergency requiring immediate ophthalmology consultation, slit-lamp examination, and targeted antimicrobial or immunosuppressive therapy.

\medterm{Intermediate Uveitis}-- Inflammation primarily affecting the vitreous body, peripheral retina, and ciliary body, with the classic presentation being pars planitis, which is characterized by the presence of snowbanking (a white, inflammatory exudate on the inferior pars plana) and snowballs (vitreous inflammatory aggregates). Patients typically present with floaters and blurred vision, often with minimal pain or redness, and slit-lamp examination reveals vitreous cells and haze, with a quiet antequiet anterior chamber. The main differential diagnosis includes multiple sclerosis, sarcoidosis, tuberculosis, Lyme disease, and cat-scratch disease, as well as the masquerade syndrome (vitreoretinal lymphoma). The diagnosis is made by slit-lamp examination showing vitreous cells and haze, with OCT used to detect cystoid macular edema, the most common cause of vision loss in intermediate uveitis.

\medterm{Intraocular Lens}-- An intraocular lens is an artificial, biocompatible, transparent lens implanted during cataract surgery to replace the natural crystalline lens that has been removed, restoring the eye's refractive power and enabling clear vision. Intraocular lenses are typically made of polymethylmethacrylate, silicone, or acrylic materials, and they are available in various designs including monofocal (correcting vision at a single distance), multifocal (providing vision at multiple distances), and toric (corrcorrecting astigmatism. The power of the IOL is calculated using biometric measurements of the eye, including axial length, corneal curvature, and anterior chamber depth, with formulas such as the SRK/T, Holladay, and Haigis formulas. The most common IOL implantation technique is in-the-bag placement within the intact capsular bag.

\medterm{Intravitreal Injection}-- An ophthalmological procedure in which medication is injected directly into the vitreous cavity of the eye through the pars plana, a safe entry zone located approximately 3.5 to 4 millimeters posterior to the limbus, to deliver therapeutic agents directly to the retina and vitreous for the treatment of various retinal diseases. The most commonly injected medications are anti-vascular endothelial growth factor agents for wet age-related macular degeneration, diabetic macmacular edema secondary to retinal vein occlusion. The procedure is performed in the clinic under sterile conditions using a topical anesthetic, and the eye is prepared with 5 percent povidone-iodine. The injection is performed through the pars plana, located 3.5 to 4 mm posterior to the limbus, using a 30-gauge or 31-gauge needle.

\medterm{Iridotomy}-- Laser peripheral iridotomy is a minimally invasive laser procedure used to treat or prevent angle-closure glaucoma by creating a small full-thickness hole in the peripheral iris, allowing aqueous humor to flow directly from the posterior chamber to the anterior chamber, bypassing the pupil and equalizing the pressure between the two chambers. The procedure is indicated for acute angle-closure glaucoma, where it is performed as an emergency to relieve pupillary block and lower intraocular pressuras an emergency to relieve pupillary block and lower intraocular pressure within minutes, and for eyes with occludable angles at high risk of angle closure. The iridotomy is created at the 10 or 2 o'clock position to avoid the visual axis and is typically 100 to 200 micrometres in diameter.

\medterm{Iritis}-- Inflammation of the iris, the most common form of anterior uveitis. Presents with acute onset eye pain, photophobia, blurred vision, and redness (ciliary flush). Signs: cells and flare in the anterior chamber, keratic precipitates (KP), synechiae (adhesions between iris and lens). Causes: idiopathic, trauma, HLA-B27-associated (ankylosing spondylitis, reactive arthritis, inflammatory bowel disease), sarcoidosis, herpes simplex/zoster. Diagnosis: slit lamp examination. Treatment: topical corticosteroids (prednisolone acetate) and cycloplegics (homatropine, atropine) to reduce inflammation and prevent synechiae.

\medterm{Juvenile angiofibroma}-- A rare, benign but locally aggressive vascular tumor of the nasopharynx, accounting for 0.05-0.5 percent of all head and neck tumors, occurring almost exclusively in adolescent males (mean age 14 years, 80-90 percent male predominance). The tumor originates from the posterolateral wall of the nasal cavity at the sphenopalatine foramen and is composed of vascular channels (irregular, thin-walled vessels lacking smooth muscle and elastic fibers) embedded in a fibrous stroma. Although histologicallhistologically benign, the tumor is locally aggressive, eroding bone and extending into the paranasal sinuses, orbit, and cranial cavity. Clinical presentation is with progressive, recurrent, typically unilateral epistaxis (the most common presenting symptom), followed by nasal obstruction, facial swelling, proptosis, and anosmia. Diagnosis is by nasal endoscopy, CT, and angiography.

\medterm{Keratitis}-- Inflammation of the cornea. Infectious keratitis: bacterial (contact lens wear, trauma: Pseudomonas, Staph aureus), viral (HSV, VZV, adenovirus), fungal (Candida, Fusarium, Aspergillus), and Acanthamoeba (contact lens-related). Presents with eye pain, photophobia, tearing, conjunctival injection, and corneal opacity/ulcer. Diagnosis: slit lamp exam with fluorescein staining reveals corneal epithelial defect. Treatment: topical antibiotics (fluoroquinolones, fortified aminoglycosides) for bacterial; antivirals (acyclovir, gancyclovir) for herpetic; antifungals for fungal. Corneal perforation requires urgent ophthalmology referral and possible penetrating keratoplasty.

\medterm{Keratoconus}-- A progressive, non-inflammatory ectatic corneal disorder characterized by bilateral (90 percent of cases) thinning and cone-like protrusion of the cornea, leading to irregular astigmatism, myopia, and decreased visual acuity. The condition typically presents in adolescence or early adulthood (peak incidence at age 15-25) and progresses until the 4th-5th decade of life, when it usually stabilizes. The prevalence is approximately 1 in 2000 in the general population. The pathophysiology involves weweakening of the corneal stroma due to reduced structural integrity of the collagen framework, with a genetic predisposition (association with Down syndrome, Ehlers-Danlos syndrome, Marfan syndrome, and mutations in the VSX1 and LOX genes). The diagnosis is made by slit-lamp examination (showing corneal thinning, Vogt striae, Fleischer ring, and apical scarring), corneal topography, and pachymetry.

\medterm{Lacrimal Apparatus}-- Produces, distributes, and drains tears for corneal moisture and protection. See \hyperlink{Anatomy}{anatomy}.

\medterm{Laser Trabeculoplasty}-- A minimally invasive laser procedure used in the treatment of open-angle glaucoma that lowers intraocular pressure by improving the outflow of aqueous humor through the trabecular meshwork. Selective laser trabeculoplasty, the current standard technique, uses a frequency-doubled Q-switched Nd:YAG laser to deliver non-thermal laser burns to the pigmented trabecular meshwork cells, stimulating biological changes that enhance aqueous humor outflow through the conventional tthe conventional outflow pathway, enhancing the outflow of aqueous humor through the trabecular meshwork without damaging adjacent tissues. SLT is most effective in primary open-angle glaucoma, pseudoexfoliation glaucoma, and pigmentary glaucoma, with minimal post-operative pain and inflammation.

\medterm{LASIK}-- LASIK, or laser-assisted in situ keratomileusis, is a refractive surgical procedure that corrects myopia, hyperopia, and astigmatism by reshaping the corneal stroma using an excimer laser, after creating a thin corneal flap with either a microkeratome (a mechanical blade) or a femtosecond laser. The procedure begins with the creation of a hinged corneal flap approximately 100 to 160 micrometers thick, which is folded back to expose the underlying corneal stroma, and the excimer laser then ablateablates the stromal tissue according to a pre-programmed treatment profile that reshapes the corneal curvature to correct the refractive error, with sub-micron precision in the removal of stromal tissue, with the depth and pattern determined by the refractive error, the optical zone, and the transition zone.

\medterm{LASIK (Laser-Assisted In Situ Keratomileusis)}-- A refractive surgical procedure that reshapes the cornea using an excimer laser to correct myopia, hyperopia, and astigmatism, thereby reducing dependence on glasses or contact lenses. A corneal flap is created with a femtosecond laser or microkeratome, the underlying stroma is ablated per preoperative topography, and the flap is repositioned. Ideal candidates are adults aged >21 with stable refraction for $\ge$1 year, adequate corneal thickness, and no active ocular disease (keratoconus, severe dry eye, autoimmune disorders). Most patients achieve 20/20 or better vision; temporary side effects include dry eye, glare, halos, and fluctuating vision. Serious complications (flap striae, infection, ectasia, epithelial ingrowth) are rare with appropriate screening and technique.

\medterm{Leber Congenital Amaurosis}-- The most severe inherited retinal dystrophy, presenting at birth or within the first few months of life with profound visual impairment, nystagmus, and pupillary light-near dissociation. Caused by mutations in over 20 genes (most commonly GUCY2D, CEP290, CRB1, RPE65, and AIPL1), encoding proteins critical for phototransduction, retinal development, or ciliary function. Inheritance is autosomal recessive. Presents with poor visual behavior (inability to fix and follow), roving nystagmus, oculodigital sign (eye poking, rubbing), and photophobia or night blindness. Fundus examination may appear normal initially, but electroretinography (ERG) is the diagnostic gold standard — shows markedly reduced or absent rod and cone responses. Treatment: gene therapy (voretigene neparvovec, Luxturna) for RPE65-associated LCA — the first FDA-approved gene therapy for a retinal disease. Low-vision rehabilitation, orientation and mobility training, and genetic counseling. Most patients have severe visual impairment from birth.

\medterm{Lisch Nodules}-- Pathognomonic ocular finding of Neurofibromatosis Type 1 (NF1, von Recklinghausen disease). Consists of clear, yellow-brown, dome-shaped melanocytic hamartomas located on the surface of the iris. Asymptomatic and do not affect vision. Present in $>90\%$ of adult NF1 patients. Part of official diagnostic criteria alongside \(\ge\)6 café-au-lait macules, neurofibromas, optic pathway glioma, and axillary/inguinal freckling (Crowe sign).

\medterm{Macula}-- The macula is a specialized region of the central retina approximately five and a half millimeters in diameter that is responsible for high-acuity central vision, color perception, and fine detail discrimination, and it contains the highest density of cone photoreceptors in the eye. The central-most region of the macula is the fovea centralis, a depression approximately one and a half millimeters in diameter where the inner retinal layers are displaced laterally, allowing light to strike the phophotoreceptors directly, maximising the efficiency of photon capture. The fovea is surrounded by the parafovea and perifovea, where the retinal layers transition to the normal architecture of the peripheral retina.

\medterm{Macular degeneration}-- A progressive degenerative disease of the central retina (macula) and the leading cause of irreversible vision loss in individuals over age 50 in developed countries, affecting approximately 10 percent of those over 65 and 30 percent over 80. The pathophysiology involves oxidative stress, chronic inflammation, and complement system dysregulation at the interface between the retinal pigment epithelium (RPE) and Bruch membrane, resulting in accaccumulation of drusen, lipofuscin, and inflammatory debris that trigger a chronic, complement-mediated immune response and progressive RPE and photoreceptor atrophy. The dry form accounts for 85-90 percent of AMD, with no effective treatment to reverse or slow the atrophy. The wet form accounts for 10-15 percent but causes the majority of severe vision loss, treated with anti-VEGF intravitreal injections.

\medterm{Miosis}-- Constriction of the pupil of the eye ($<2\text{ mm}$ diameter). Controlled by parasympathetic nerve fibers via the oculomotor nerve (CN III) acting on the sphincter pupillae muscle. Causes: Opioid toxicity (pathognomonic pinpoint pupils), organophosphate/nerve agent poisoning (cholinergic excess), Horner syndrome (loss of sympathetic input: ptosis, miosis, anhidrosis), pontine hemorrhage, and miotic eye drops (pilocarpine).

\medterm{Mydriasis}-- Dilation of the pupil of the eye ($>5\text{ mm}$ diameter). Controlled by sympathetic nerve fibers acting on the dilator pupillae muscle. Causes: Anticholinergic toxicity (atropine, scopolamine, diphenhydramine — "blind as a bat, mad as a hatter"), sympathomimetic agents (cocaine, amphetamines), third cranial nerve (CN III) palsy (uncal herniation causing "blown pupil"), and traumatic iridoplegia.

\medterm{Myopia}-- Myopia, commonly known as nearsightedness, is a refractive error in which parallel light rays from distant objects are focused anterior to the retina rather than directly on it, resulting in blurred distance vision while near vision remains relatively clear. The condition is caused by an excessive axial length of the globe relative to the refractive power of the cornea and lens, or by excessive corneal curvature, and it is corrected with concave (minus-powered) diverging lenses that move the focfocal point of light from in front of the retina back onto the retina, enabling clear distance vision. The degree of myopia is quantified in dioptres (D), with low (less than -3.00 D), moderate (-3.00 D to -6.00 D), and high (greater than -6.00 D) categories, with high myopia being associated with increased risk of retinal detachment, choroidal neovascularization, macular atrophy, glaucoma, and cataract.

\medterm{Night Blindness}-- Impaired vision in dim light (nyctalopia). Causes: vitamin A deficiency (most common cause worldwide, leading to retinol deficiency and rod dysfunction), retinitis pigmentosa (progressive rod-cone degeneration, genetic), congenital stationary night blindness, advanced glaucoma, and cataract. Vitamin A deficiency: also causes xerophthalmia (dry cornea) and Bitot spots (foamy conjunctival deposits); treat with vitamin A supplementation. Retinitis pigmentosa: no cure; vitamin A, lutein supplementation, and low-vision aids.

\medterm{Blindness}-- Severe or complete loss of vision in one or both eyes, caused by disease or injury affecting the cornea, lens, retina, optic nerve, visual pathways, or visual cortex. It may be congenital or acquired, sudden or progressive, and reversible or permanent; sudden visual loss is an ophthalmic emergency.

\medterm{Non-Contact Tonometry}-- Non-contact tonometry, commonly known as the air puff test, is a method of measuring intraocular pressure that uses a calibrated air pulse to applanate the cornea, with the force of the air pulse determined by a sensor that detects the moment of applanation, providing a rapid, non-invasive measurement without requiring direct corneal contact. The technique eliminates the need for topical anesthesia and reduces the risk of cross-infection between patients, making it particularly useful for screenscreening in high-volume settings and for patients unable to tolerate contact tonometry. The device measures the force of the air puff at the moment of corneal applanation, but it is generally less accurate than Goldmann applanation tonometry, especially in patients with irregular corneas.

\medterm{Non-Proliferative Diabetic Retinopathy}-- The early stage of diabetic retinopathy, characterized by structural abnormalities of the retinal vasculature resulting from chronic hyperglycemia-induced damage to the endothelial cells and pericytes of retinal capillaries, without neovascularization. The clinical features include microaneurysms, which are the earliest clinically visible sign and appear as small, red dots on fundus examination; dot and blot hemorrhages, representing deeper retinal hemorhemorrhages, representing deeper retinal hemorrhages; hard exudates, which are lipid deposits in the outer plexiform layer; cotton-wool spots, representing nerve fibre layer infarcts; and retinal vascular changes, including venous beading, arteriolar narrowing, and intraretinal microvascular abnormalities.

\medterm{Nuclear Sclerosis}-- The most common type of age-related cataract, characterized by progressive yellowing and hardening of the lens nucleus due to the accumulation and aggregation of insoluble crystallin proteins, resulting in a gradual decrease in vision that is often noticed as increased difficulty with distance vision and a phenomenon called second-sight vision, where temporarily improved near vision occurs as the lens becomes more myopic from its increased refractive index. The lens nucleus becomes progressively harder and more yellow, and the resulting increase in the refractive index of the nucleus causes the lens power to increase, leading to a shift toward myopia. The diagnosis is made by slit-lamp examination after dilation, showing the yellowing and increased density of the lens nucleus, graded from 1+ to 4+.

\medterm{Nystagmus}-- A rhythmic, involuntary oscillation of the eyes that can be horizontal, vertical, rotational, or mixed in pattern, and it is classified based on the underlying cause as infantile (developing in the first six months of life), acquired (developing later in life), or periodic alternating. Infantile nystagmus is the most common form, often associated with albinism, congenital cataracts, or other conditions that impair visual development, and it typically has a null point (a position of the specific direction of gaze, or a head turn or tilt that reduces the intensity of the nystagmus). In acquired nystagmus, the onset is typically subacute and oscillopsia is prominent, and neurological evaluation with imaging is essential to rule out an underlying structural or metabolic cause, including vestibular disorders, brainstem or cerebellar lesions, and drug-induced causes.

\medterm{Opacity}-- Lack of transparency of a normally transparent ocular structure (cornea, lens, vitreous, or aqueous humour). Corneal opacity results from scarring (trauma, infection—herpes simplex keratitis, trachoma), dystrophies (Fuchs endothelial dystrophy, granular/lattice/macular dystrophies), or edema. Lenticular opacity is cataract—age-related, congenital, traumatic, or metabolic (diabetes). Vitreous opacity arises from hemorrhage (diabetic retinopathy, retinal tear), inflammation (vitritis, posterior uveitis), or asteroid hyalosis (calcium-lipid deposits in the vitreous, benign). Assessment: slit-lamp examination, visual acuity, and optical coherence tomography. Opacity causes decreased visual acuity, glare, and contrast sensitivity loss. Management depends on cause: corneal transplant (penetrating keratoplasty, DSEK) for irreversible corneal opacity; cataract surgery (phacoemulsification with intraocular lens implantation) for visually significant cataract; vitrectomy for non-clearing vitreous hemorrhage.

\medterm{Optic Atrophy}-- End-stage degeneration of the optic nerve fibers resulting from any optic neuropathy. Causes: glaucoma, ischemic optic neuropathy, multiple sclerosis, compressive lesions (pituitary tumor, aneurysm), toxins (methanol, ethambutol), and hereditary (Leber hereditary optic neuropathy). Fundoscopic findings: pale optic disc with sharp margins (primary atrophy) or blurred margins (secondary atrophy). Presents with visual acuity loss, color vision deficits, and visual field defects. Irreversible; treatment addresses underlying cause.

\medterm{Optic Disc}-- The point on the retina where the optic nerve exits the eye, lacking photoreceptors and creating the natural blind spot.

\medterm{Optic Neuritis}-- An inflammatory demyelinating condition of the optic nerve that typically presents with acute or subacute unilateral vision loss, pain with eye movement, and a relative afferent pupillary defect, and it is most commonly associated with multiple sclerosis, occurring in approximately fifty percent of patients at some point during their disease course. The pathophysiology involves autoimmune-mediated demyelination of the optic nerve, with T lymphocyte infiltration and antibody-medantibody-mediated demyelination. The afferent visual dysfunction is characterized by reduced best-corrected visual acuity, decreased contrast sensitivity, dyschromatopsia (impaired color vision, particularly red color desaturation), and a relative afferent pupillary defect. A central or paracentral scotoma is found on formal visual field testing.

\medterm{Optical Coherence Tomography}-- A non-invasive imaging technique that uses low-coherence interferometry to produce high-resolution, cross-sectional images of the retina and other ocular structures, providing quantitative measurements of retinal layer thickness and detecting structural abnormalities with micrometer-level resolution. In ophthalmology, OCT is the gold standard for diagnosing and monitoring macular diseases including age-related macular degeneration, diabetic macular edema, and macumacular hole, epiretinal membrane, and central serous chorioretinopathy, and for measuring retinal nerve fibre layer thickness in glaucoma. The principle of OCT is analogous to ultrasound imaging but using light rather than sound, achieving an axial resolution of 2 to 7 micrometres.

\medterm{Optometry}-- A healthcare profession concerned with examining the eyes and visual system for defects and abnormalities, prescribing corrective lenses (spectacles, contact lenses), diagnosing eye disease, and providing pre- and post-operative care for surgical eye treatments. Optometrists (in the UK) perform sight tests (refraction, visual field assessment, tonometry, retinal examination), detect ocular disease (glaucoma, diabetic retinopathy, macular degeneration, cataracts), and refer to ophthalmologists when medical or surgical intervention is required. They also manage low vision, binocular vision disorders (strabismus, amblyopia), and provide advice on visual ergonomics. In the UK, optometrists train for 4 years (BSc Optometry) followed by a pre-registration year and professional qualifying examinations.

\medterm{Orbital Cellulitis}-- A serious infection of the orbital tissues posterior to the orbital septum, potentially sight-threatening and life-threatening due to the risk of optic nerve compression, cavernous sinus thrombosis, and intracranial extension. The condition is classified by the Chandler system into five stages: preseptal cellulitis (infection confined to the eyelid and tissues anterior to the orbital septum), orbital cellulitis (infection of the orbital fat and muscles without abscess formaformation), and postseptal cellulitis (the true orbital cellulitis, involving the orbital fat, muscles, optic nerve, and neurovascular structures). Preseptal cellulitis presents with eyelid erythema, edema, warmth, and tenderness, with no proptosis, no limitation of extraocular movements, no relative afferent pupillary defect, and no visual loss, and it is managed with oral antibiotics.

\medterm{Palpebral Conjunctiva}-- The palpebral conjunctiva is the portion of the conjunctival membrane that lines the inner surface of the upper and lower eyelids, and it is subdivided into the tarsal conjunctiva, which is firmly adherent to the tarsal plates and appears smooth and pink, and the forniceal conjunctiva, which forms loose, redundant folds at the superior and inferior fornices where the palpebral conjunctiva reflects to become the bulbar conjunctiva. The palpebral conjunctiva contains meibomian glands (sebaceous glmeibomian glands that produce the lipid layer of the tear film. The tarsal conjunctiva is normally pink and smooth, and abnormalities such as follicles (characteristically seen in viral and chlamydial conjunctivitis), papillae (characteristically seen in allergic conjunctivitis), and membranes provide important diagnostic clues.

\medterm{Panretinal Photocoagulation}-- A laser treatment for proliferative diabetic retinopathy that involves applying laser burns to the peripheral retina (typically 1,200 to 2,000 burns over three to four sessions) to ablate ischemic retinal tissue, reducing the production of vascular endothelial growth factor and other angiogenic factors that drive pathological neovascularization. The laser burns destroy the oxygen-consuming photoreceptor cells in the peripheral retina, improving oxygenation of the rremaining retina, reducing the VEGF drive for neovascularization, thereby causing regression of existing new vessels and preventing the growth of new abnormal vessels. The specific indications for PRP are high-risk proliferative diabetic retinopathy.

\medterm{Panuveitis}-- Inflammation involving all segments of the uveal tract, including the anterior chamber, vitreous cavity, and choroid, resulting in a severe, potentially sight-threatening condition that requires aggressive systemic immunosuppressive therapy. The clinical presentation combines features of anterior, intermediate, and posterior uveitis, with patients experiencing significant eye pain, photophobia, blurred vision, and decreased vision, and examination revealing cells and flare in the aanterior chamber, vitreous haze and cells, and active choroiditis or chorioretinitis, often with exudative retinal detachment, cystoid macular edema, and papillitis. The etiology includes infectious causes (toxoplasmosis, tuberculosis, syphilis, cytomegalovirus), non-infectious causes (sarcoidosis, Behcet disease, Vogt-Koyanagi-Harada syndrome, birdshot chorioretinopathy), and masquerade syndrome (intraocular lymphoma).

\medterm{Papilledema}-- Swelling of the optic disc caused by elevated intracranial pressure, and it is a critical clinical sign that requires urgent investigation to identify the underlying cause, which may be life-threatening. The pathophysiology involves transmission of elevated cerebrospinal fluid pressure through the optic nerve sheath to the optic nerve head, causing axoplasmic stasis, disc edema, and potentially progressive optic nerve damage. Clinical findings include disc blurring, elevation, venvenous congestion, retinal venous pulsation loss, and peripapillary haemorrhages. The urgency mandates immediate neuroimaging (CT head without contrast) to rule out a space-occupying lesion or hydrocephalus, followed by lumbar puncture with opening pressure measurement if CT is negative.

\medterm{Phacoemulsification}-- The most commonly performed modern cataract surgery technique, involving the use of an ultrasonic handpiece with a titanium tip that vibrates at high frequency to fragment and emulsify the opaque lens nucleus while simultaneously irrigating and aspirating the lens material through a coaxial or bimanual system, leaving the posterior lens capsule intact for implantation of an artificial intraocular lens. The procedure is performed through a small self-sealing incision, typictypically 2.2 to 3.2 mm in width, through the clear cornea near the limbus. Phacoemulsification has a low rate of wound leak, post-operative inflammation, and induced astigmatism, with rapid visual recovery and a return to full activity within days to weeks.

\medterm{Pinguecula}-- A yellowish-white, slightly elevated benign degenerative lesion of the bulbar conjunctiva near the limbus, from elastotic degeneration of the substantia propria. Chronic ultraviolet exposure, dust, wind, and dryness are the principal risk factors. Most are asymptomatic but may cause irritation, foreign body sensation, cosmetic concern, or inflammation (pingueculitis). Management is conservative with lubrication and UV protection; excision is rarely required.

\medterm{Pink Eye}-- Inflammation of the conjunctiva producing a pink or red appearance of the eye, most often indicating conjunctivitis. Causes are infectious (viral, bacterial, chlamydial) or noninfectious (allergic, irritant), with viral the most common overall. Typical features include conjunctival injection, discharge, tearing, and foreign body sensation. Most cases are self-limited or managed with targeted therapy and hygiene to limit spread.

\medterm{Posterior Subcapsular Cataract}-- A type of lens opacity located at the posterior aspect of the lens, just anterior to the posterior lens capsule, in the visual axis, and it is particularly visually disabling because of its central location. The opacity appears as a granular, plaque-like whitish area on the posterior lens surface on slit-lamp examination, and it may have a flower-petal or vacuolar pattern. This form of cataract is more common in younger patients and is strongly associated with ccorticosteroid use (both topical, inhaled, and systemic, which is the strongest modifiable risk factor, with risk increasing with dose and duration of use), diabetes mellitus, trauma, intraocular inflammation, and radiation exposure. The treatment is surgical removal with phacoemulsification and IOL implantation.

\medterm{Posterior Uveitis}-- Inflammation primarily affecting the choroid (choroiditis), retina (retinitis), or both (chorioretinitis), and it can be caused by infectious agents including Toxoplasma gondii, cytomegalovirus, herpes simplex virus, tuberculosis, and syphilis, as well as non-infectious causes including sarcoidosis, Behcet disease, and Vogt-Koyanagi-Harada syndrome. The clinical presentation depends on the underlying cause and may include floaters, blurred vision, scotomas, and decreased visdecreased vision with a relative afferent pupillary defect in severe or unilateral cases. The diagnosis is based on fundoscopic examination showing active choroidal or retinal lesions, with ancillary testing to characterise the disease and laboratory evaluation to identify the underlying cause.

\medterm{Preseptal Cellulitis}-- Preseptal cellulitis, also known as periorbital cellulitis, is an infection of the eyelid and periorbital tissues anterior to the orbital septum, without involvement of the orbital contents, and it is the most common orbital infection, particularly in children. The condition typically results from contiguous spread from adjacent skin infections, insect bites, or trauma, and common causative organisms include Staphylococcus aureus and Streptococcus pyogenes in community-acquired cases, and HaemopHaemophilus influenzae in children and Staphylococcus aureus and Streptococcus pyogenes in adults. Preseptal cellulitis presents with eyelid edema, erythema, warmth, and tenderness, with no proptosis, no limitation of extraocular movements, no relative afferent pupillary defect, and no visual loss. Management is with oral antibiotics covering the common pathogens.

\medterm{PRK}-- Photorefractive keratectomy, commonly known as PRK, is a refractive surgical procedure that corrects myopia, hyperopia, and astigmatism by reshaping the corneal stroma using an excimer laser after mechanical removal of the corneal epithelium, without creating a corneal flap. The procedure begins with the removal of the corneal epithelium using an alcohol solution, mechanical debridement, or the excimer laser itself, and the excimer laser then ablates the exposed stromal surface according to a coa customized laser treatment profile, removing the Bowman layer and anterior stroma without creating a flap. After ablation, a bandage contact lens is placed to promote epithelial healing. The epithelial defect typically heals over 3 to 5 days, and visual recovery is more gradual than LASIK.

\medterm{Proliferative Diabetic Retinopathy}-- The advanced stage of diabetic retinopathy, characterized by the growth of new, fragile blood vessels (neovascularization) on the retinal surface, optic disc, or iris in response to severe retinal ischemia and the release of vascular endothelial growth factor. These new vessels are abnormal in structure, lacking tight junctions and pericyte support, making them prone to hemorrhage into the vitreous cavity, which can cause sudden, painless, severe vision lossloss due to vitreous hemorrhage. The neovascular vessels also grow on the iris (rubeosis iridis) and in the anterior chamber angle, leading to neovascular glaucoma, a severe, refractory form of glaucoma requiring prompt panretinal photocoagulation to induce regression of the new vessels.

\medterm{Prostaglandin Analogs}-- A class of topical medications used in the treatment of open-angle glaucoma and ocular hypertension that lower intraocular pressure by increasing the uveoscleral outflow of aqueous humor, reducing intraocular pressure by approximately twenty-five to thirty-five percent. The mechanism involves binding to prostaglandin F2-alpha receptors in the ciliary muscle and ciliary body, which activates matrix metalloproteinases that remodel the extracellular matrix of the uveosclerthe uveoscleral pathway through enhanced matrix metalloproteinase activity, reducing the outflow resistance. Available agents include latanoprost (the most commonly prescribed), travoprost, tafluprost, and bimatoprost, and they are generally safe and well tolerated, with the most common side effect being conjunctival hyperaemia.

\medterm{Pterygium}-- A triangular, fibrovascular proliferation of bulbar conjunctiva extending onto the cornea, usually nasal, associated with chronic ultraviolet exposure, dust, and dry climates. Most are asymptomatic or cause mild irritation, but growth toward the visual axis can induce astigmatism and obstruct vision. Excision with conjunctival autograft is reserved for threat to the visual axis or significant astigmatism.

\medterm{Ptosis}-- The drooping of the upper eyelid, which can be congenital or acquired, and can range from mild drooping that is primarily a cosmetic concern to severe drooping that obstructs the visual axis and causes amblyopia in children or visual impairment in adults. Congenital ptosis is most commonly caused by developmental abnormalities of the levator palpebrae superioris muscle, including myogenic dysgenesis, and is typically present at birth with poor levator function and an absent lid crease.congenital ptosis is most commonly caused by developmental abnormalities of the levator palpebrae superioris muscle, including myogenic dysgenesis, and is typically present at birth with poor levator function and an absent lid crease. Acquired ptosis is further classified by etiology into aponeurotic (most common in older adults, caused by dehiscence of the levator aponeurosis), myogenic, neurogenic, and mechanical ptosis.

\medterm{Pupil}-- Central round iris aperture for light entry, size controlled by sphincter pupillae (CN III parasympathetic, constricts to 2-3 mm) and dilator pupillae (sympathetic, dilates to 8 mm). See \hyperlink{Anatomy}{anatomy}.

\medterm{Refraction}-- The bending of light rays as they pass through the eye's optical system (cornea, aqueous humour, crystalline lens, vitreous humour) to focus a sharp image on the retina. The refractive power of the eye is measured in dioptres (D). Emmetropia: normal refractive state—parallel light rays focus exactly on the retina without accommodation. Ametropia (refractive error): (1) Myopia (short-sightedness)—light focuses in front of the retina, corrected by concave (negative) lenses; axial myopia (elongated eyeball) is most common. (2) Hypermetropia (long-sightedness)—light focuses behind the retina, corrected by convex (positive) lenses; young hypermetropes may compensate by accommodation. (3) Astigmatism—unequal curvature of the cornea or lens causes light to focus on multiple points, corrected by cylindrical lenses. (4) Presbyopia—age-related loss of accommodation (lens hardening), typically becomes symptomatic after age 40, corrected by reading glasses or bifocals. Refraction is measured objectively (retinoscopy, autorefractor) and subjectively (the patient's response to lens choices during a sight test). Refractive surgery (LASIK, PRK, SMILE) reshapes the cornea to correct refractive errors.

\medterm{Refractive Errors and Astigmatism}-- Refractive errors are vision disorders caused by the eye's inability to properly focus light onto the retina, resulting in blurred vision. Common types include myopia (nearsightedness), hyperopia (farsightedness), presbyopia (age-related loss of near vision), and astigmatism, which occurs when the cornea or lens has an irregular curvature that causes distorted vision at all distances. These conditions are typically diagnosed through a comprehensive eye examination and are correctable with eyeglasses, contact lenses, or refractive surgery such as LASIK.

\medterm{Refractive Surgery}-- Surgical procedures that reshape the cornea or alter the eye to correct refractive errors and reduce dependence on eyeglasses or contact lenses. Laser techniques include LASIK (corneal flap with excimer ablation), photorefractive keratectomy (surface ablation), and small-incision lenticule extraction (SMILE). Phakic intraocular lens implantation suits patients unsuited to corneal laser surgery. Careful candidate selection avoids ectasia, infection, and dry eye.

\medterm{Retinal Detachment}-- A separation of the neurosensory retina from the underlying retinal pigment epithelium, creating a potential space that fills with fluid and resulting in progressive vision loss in the affected area, and it is classified into three main types based on the underlying mechanism. Rhegmatogenous retinal detachment, the most common type, is caused by a full-thickness break or tear in the retina that allows liquefied vitreous to enter the subretinal space, separating the retina ffrom the underlying RPE. Tractional retinal detachment is caused by the contraction of fibrovascular membranes on the retinal surface, most commonly in proliferative diabetic retinopathy. Exudative retinal detachment occurs when subretinal fluid accumulates due to leakage from blood vessels, without a retinal break, most commonly associated with inflammatory conditions.

\medterm{Retinal Layers}-- The retina consists of ten distinct layers that process visual information through a complex neural circuit. From innermost to outermost, these layers include the internal limiting membrane (the boundary between retina and vitreous), nerve fiber layer (containing retinal ganglion cell axons), ganglion cell layer (containing retinal ganglion cell bodies), inner plexiform layer (synapses between bipolar and ganglion cells), inner nuclear layer (cell bodies of bipolar, horizontal, amacrine, and MllMuller cells), outer plexiform layer (synapses between photoreceptors and bipolar cells), outer nuclear layer (cell bodies of rod and cone photoreceptors), external limiting membrane, and the photoreceptor layer. The retinal pigment epithelium (RPE) is a single layer of pigmented cells beneath the photoreceptors that is essential for photoreceptor health and function.

\medterm{Retinitis Pigmentosa}-- A group of inherited retinal dystrophies characterized by progressive degeneration of rod and cone photoreceptors, leading to visual field loss and eventual blindness. Prevalence: 1 in 4,000. Genetics: autosomal dominant, autosomal recessive, X-linked, and mitochondrial inheritance; >60 genes implicated (RHO, RPGR, RP2, USH2A, PDE6B, ABCA4). Pathophysiology: rod photoreceptors degenerate first, followed by cones. Characteristic fundus findings: 'bone-spicule' pigmentary deposits in the mid-peripheral retina (perivascular), waxy pallor of the optic disc, and attenuated retinal arterioles. Clinical: night blindness (nyctalopia—earliest symptom, often in childhood), progressive constriction of the visual field (tubular vision—preserves central vision until late stages), and eventually loss of central vision (legal blindness by age 40--60). Associated systemic conditions: Usher syndrome (RP plus congenital sensorineural hearing loss), Bardet–Biedl syndrome (RP, obesity, polydactyly, intellectual disability, renal abnormalities). Diagnosis: electroretinography (ERG—reduced or extinguished rod and cone responses), visual field testing (Goldmann perimetry—constricted fields), optical coherence tomography (OCT—thinning of outer retinal layers), and genetic testing. Treatment: no cure. Management: low vision aids, vitamin A palmitate (15,000 IU/day—may slow progression in some forms, controversial), and avoidance of bright light. Gene therapy: voretigene neparvovec (Luxturna) approved for RPE65-related RP (restores visual function). Retinal prostheses (Argus II—epiretinal implant) for end-stage disease. Research: stem cell therapy, optogenetics. Prognosis: progressive, most patients are legally blind by age 60.

\medterm{Rhegmatogenous Retinal Detachment}-- The most common type of retinal detachment, caused by a full-thickness break or tear in the retina that allows liquefied vitreous humor to enter the subretinal space, separating the neurosensory retina from the underlying retinal pigment epithelium. The condition typically occurs in association with posterior vitreous detachment, where vitreoretinal traction at sites of abnormal adhesion causes retinal tears, most commonly at the equator of the eye. Patients typically present with sudden onset of unilateral, painless, progressive visual field loss, often described as a curtain or shade over their vision, with or without preceding photopsias and floaters. On examination, there is a relative afferent pupillary defect, decreased visual acuity if the macula is detached, and the neurosensory retina is seen to be elevated and billowing on fundoscopy.

\medterm{Scotoma}-- A focal area of diminished or absent vision within an otherwise intact visual field, classified by location and morphology. Central scotomas involving fixation are characteristic of macular disease and optic neuropathies; paracentral and arcuate scotomas are typical of glaucoma; peripheral scotomas result from retinal detachment, chorioretinitis, or visual pathway lesions. Perimetry localizes the site of damage.

\medterm{Slit Lamp Biomicroscopy}-- A fundamental ophthalmological examination technique that uses a high-intensity, adjustable slit beam of light combined with a binocular microscope to provide a magnified, stereoscopic view of the anterior segment and, with the use of auxiliary lenses, the posterior segment of the eye. The slit beam can be adjusted in width, height, and angle, allowing the examiner to create an optical section through the transparent ocular media and visualize individual corneal layersand structures in fine detail, including the corneal epithelium, Bowman layer, stroma, Descemet membrane, and endothelium, as well as allowing the detection of cells and flare in the anterior chamber, cataract in the lens, and vitreous cells. The slit lamp is also used for Goldmann applanation tonometry and for performing diagnostic and therapeutic procedures.

\medterm{SMILE}-- Small incision lenticule extraction, commonly known as SMILE, is a minimally invasive refractive surgical procedure that corrects myopia and myopic astigmatism by creating and removing a thin, disc-shaped piece of corneal stroma called a lenticule through a small incision, without creating a corneal flap. The procedure uses a femtosecond laser to create two intrastromal cuts: the posterior surface of the lenticule and the anterior cap, and a small peripheral incision approximately 2 to 4 millimemillimetres in length creating a side-entry channel for extraction of the lenticule. SMILE offers advantages over LASIK, including the absence of a corneal flap, preservation of corneal biomechanical stability, and a lower risk of dry eye syndrome.

\medterm{Squint}-- A misalignment of the visual axes of the two eyes, where one eye is deviated inward (esotropia), outward (exotropia), upward (hypertropia), or downward (hypotropia) relative to the other, disrupting binocular vision and potentially leading to amblyopia in children. See \hyperlink{Strabismus}{Strabismus}.

\medterm{Strabismus}-- A misalignment of the visual axes of the two eyes, where one eye is deviated inward (esotropia), outward (exotropia), upward (hypertropia), or downward (hypotropia) relative to the other, disrupting binocular vision and potentially leading to amblyopia in children. The condition can be constant or intermittent, unilateral or alternating, and it may be comitant (the angle of deviation is the same in all directions of gaze) or incomitant (the angle varies depending on the direction odirection of gaze, indicating a paretic or restrictive cause). The diagnosis is made by the cover-uncover test and the alternate cover test, which detect manifest or latent strabismus and measure the angle of deviation in prism dioptres. Additional tests include measurement of visual acuity, evaluation of ductions and versions, and the Hess chart.

\medterm{Stye}-- A common infection of the eyelid. External hordeolum: acute bacterial infection of the Zeis or Moll glands at the eyelash follicle, typically caused by Staphylococcus aureus. Presents as a tender, red, well-localized swelling at the lid margin. Internal hordeolum: infection of the meibomian gland, larger and more painful than external. Treatment: warm compresses 3-4 times daily, lid hygiene. Incision and drainage for persistent cases. Avoid squeezing to prevent spread of infection.

\medterm{Trabeculectomy}-- A conventional glaucoma filtration surgery that creates a new drainage pathway for aqueous humor by removing a small piece of the trabecular meshwork and creating a partial-thickness scleral flap under which aqueous humor filters from the anterior chamber into the subconjunctival space, forming a filtering bleb. The procedure is indicated for glaucoma that is uncontrolled with maximal medical and laser therapy, and it achieves intraocular pressure reduction to target levels in target levels in up to 80 to 90 percent of patients at 2 years. The procedure involves creating a conjunctival flap, a partial-thickness scleral flap, excising a trabecular and deep scleral block, and creating a peripheral iridectomy. The scleral flap is then sutured to control the rate of aqueous filtration.

\medterm{Uvea}-- The middle vascular coat of the eye, consisting of three parts from anterior to posterior: the iris (the coloured part of the eye, controlling pupil size), the ciliary body (producing aqueous humour and controlling lens shape via the zonular fibers), and the choroid (the vascular layer supplying the outer retina, between the sclera and the retinal pigment epithelium). The uvea is the site of inflammation in uveitis. Uveal melanoma is the most common primary intraocular malignancy in adults, arising from melanocytes in the choroid, ciliary body, or iris. Benign uveal tumors include naevi.

\medterm{Uveitis}-- An inflammatory condition of the uveal tract, which includes the iris, ciliary body, and choroid, and it is classified anatomically based on the location of the inflammation: anterior uveitis (iritis and iridocyclitis), intermediate uveitis (pars planitis and periphlebitis), posterior uveitis (choroiditis and chorioretinitis), and panuveitis (involvement of all three segments). Anterior uveitis is the most common form, presenting with eye pain, photophobia, blurred vision, and rednessredness with ciliary flush. The diagnosis is confirmed by slit-lamp examination revealing cells and flare in the anterior chamber, keratic precipitates, and synechiae. Treatment depends on location and severity, with topical corticosteroids for anterior uveitis and systemic immunosuppression for posterior and panuveitis.

\medterm{Viral Conjunctivitis}-- The most common form of infectious conjunctivitis, caused predominantly by adenoviruses (serotypes 3, 4, 7, 8, 19, and 37), and it is highly contagious, spreading through direct contact with infected ocular secretions or contaminated surfaces. The condition typically begins unilaterally and spreads to the contralateral eye within one to two days, presenting with watery discharge, foreign body sensation, tearing, and conjunctival injection, often accompanied by preauriculapreauricular lymphadenopathy in 30 to 50 percent of cases. The diagnosis is clinical, and treatment is largely supportive (cool compresses, lubricating drops, and careful hand hygiene to prevent spread), as the disease is self-limited and resolves within 2 to 3 weeks.

\medterm{Vision}-- The sensory ability to perceive light and interpret visual information processed by the eyes and brain, enabling recognition of shapes, colors, depth, and motion. Normal vision depends on the proper function of the cornea, lens, retina, and optic nerve, as well as intact neural pathways to the visual cortex. Age-related changes such as presbyopia and cataracts commonly affect vision, while conditions like glaucoma, macular degeneration, and diabetic retinopathy can cause permanent vision loss if not detected and treated early.

\medterm{Vision and Hearing}-- The two primary sensory modalities through which humans perceive and interact with their environment. Age-related declines in both senses are common, with presbyopia and cataracts affecting vision and presbycusis affecting hearing, often beginning as early as the fourth decade of life. Regular sensory screening is important for early detection of impairments, and corrective measures including eyeglasses, contact lenses, hearing aids, and cochlear implants can substantially improve quality of life and reduce the risk of falls, social isolation, and cognitive decline.

\medterm{Vision Loss}-- A decrease in the ability to see that cannot be fully corrected with eyeglasses, contact lenses, or medical treatment. It ranges from partial impairment to complete blindness and can result from conditions such as macular degeneration, glaucoma, diabetic retinopathy, cataracts, and traumatic injury. Low vision rehabilitation services, assistive technologies including magnifiers and screen readers, and environmental modifications help individuals with vision loss maintain independence and quality of life.

\medterm{Visual Acuity}-- The resolving power of the eye, defined as the ability to distinguish two points as distinct, and the most common clinical measure of visual function. It is quantified using the Snellen chart (e.g., 20/20), in which the numerator is the testing distance and the denominator the distance a normal eye could read the same line; smaller denominators indicate better acuity.

\medterm{Visual Field Defects}-- Areas of reduced or absent vision within the visual field that correspond to damage at specific points along the visual pathway, and their pattern provides crucial localizing information for diagnosis. A scotoma is an area of diminished or absent vision surrounded by normal vision, with the most common types being central scotomas (from macular disease or optic neuropathy) and paracentral scotomas (from glaucoma or retinal disease). Homonymous hemianopia is loss of the sthe same side of the visual field in both eyes, which localizes to the optic tract or optic radiation contralateral to the side of the field loss, most commonly caused by stroke. Bitemporal hemianopia is loss of the temporal half of the visual field in both eyes, which localizes to the optic chiasm, most commonly caused by a pituitary adenoma.

\medterm{Visual Field Perimetry}-- A systematic method of measuring the sensitivity of the retina across the entire visual field using standardized light stimuli, providing a quantitative map of visual function that can detect, characterize, and monitor visual field loss caused by diseases of the retina, optic nerve, and visual pathways. Automated static perimetry, most commonly performed using the Humphrey Field Analyzer, presents points of light at varying intensities throughout the visual field and rereports the threshold sensitivity at each tested point. The Humphrey 24-2 and 30-2 protocols are the most commonly used. The results are displayed as a gray-scale map, a total deviation probability map, and a pattern deviation probability map.

\medterm{Vitrectomy}-- Surgical removal of the vitreous gel, typically via the pars plana approach, using a cutting probe under microscopic visualization. Indications include vitreous hemorrhage, complicated retinal detachment, macular holes, epiretinal membranes, intraocular foreign bodies, and endophthalmitis. The three-port trocar technique often employs endolaser and gas or silicone oil tamponade. Complications include retinal tears, cataract, and elevated intraocular pressure.

\medterm{Vitreous Body}-- The vitreous humor, or vitreous body, is a transparent, gel-like substance that fills the space between the lens and the retina, occupying approximately eighty percent of the volume of the eye, and it consists primarily of water (approximately ninety-nine percent) with a network of collagen fibrils and hyaluronic acid that give it its gel-like consistency. See \hyperlink{Vitreous Humor}{Vitreous Humor}.

\medterm{Vitreous Humor}-- The vitreous humor, or vitreous body, is a transparent, gel-like substance that fills the space between the lens and the retina, occupying approximately eighty percent of the volume of the eye, and it consists primarily of water (approximately ninety-nine percent) with a network of collagen fibrils and hyaluronic acid that give it its gel-like consistency. The vitreous is attached to the retina at several sites, including the vitreous base anteriorly (the strongest attachment), the optic disc mamargin (the ring of attachment around the optic disc). The vitreous is subject to age-related changes, with liquefaction (synchysis senilis) and posterior vitreous detachment, which can be asymptomatic or cause flashes and floaters.

\medterm{Xerophthalmia}-- See clinical\_nutrition.

\medterm{YAG Capsulotomy}-- YAG laser posterior capsulotomy is a minimally invasive laser procedure used to treat posterior capsule opacification, the most common late complication of cataract surgery, in which the posterior lens capsule becomes opacified by proliferation of residual lens epithelial cells, causing decreased vision. The procedure uses a neodymium:yttrium-aluminum-garnet laser to create a central opening in the opacified posterior capsule, allowing light to pass unobstructed to the retina and restoring clear vision, often within 24 to 48 hours. The procedure is performed on an outpatient basis at the slit lamp, with the pupil dilated and topical anesthetic applied. The laser creates an opening in the posterior capsule typically measuring 3 to 5 mm in diameter, with a total of 20 to 50 laser applications of 1 to 5 mJ each.

\medterm{Intraocular Pressure (IOP)}-- The pressure generated by aqueous humor within the eye, normally maintained by a balance between aqueous production and outflow. Abnormal IOP is an important risk factor for glaucomatous optic neuropathy, but glaucoma can occur at statistically normal pressures.

\medterm{Ocular Hypertension}-- Persistently elevated intraocular pressure without demonstrable glaucomatous optic-nerve damage or corresponding visual-field loss. It increases future glaucoma risk and requires individualized assessment of corneal thickness, optic-nerve appearance, pressure, and other risk factors.

\medterm{Primary Open-Angle Glaucoma}-- A chronic optic neuropathy with characteristic retinal ganglion-cell and optic-nerve loss, usually accompanied by progressive visual-field defects, despite an anatomically open anterior-chamber angle. Intraocular pressure is an important modifiable risk factor, and treatment lowers pressure to slow progression.

\medterm{Acute Angle-Closure Glaucoma}-- An ophthalmic emergency in which the anterior-chamber angle closes and aqueous outflow is abruptly obstructed. Severe ocular pain, headache, nausea, blurred vision, halos, corneal edema, a fixed mid-dilated pupil, and markedly elevated IOP may occur; urgent pressure-lowering therapy and definitive iridotomy are required.

\medterm{Relative Afferent Pupillary Defect (RAPD)}-- An asymmetric reduction in the direct pupillary light response caused by unequal afferent input from the retina or optic nerve. It is detected with the swinging-flashlight test and is also called a Marcus Gunn pupil.

\medterm{Optic Disc Edema}-- Swelling of the optic-nerve head caused by axoplasmic-flow disruption. Bilateral optic-disc edema from raised intracranial pressure is papilledema; unilateral or non-pressure causes require separate evaluation for optic-nerve, retinal, inflammatory, ischemic, or infiltrative disease.

\medterm{Posterior Vitreous Detachment (PVD)}-- Separation of the posterior vitreous cortex from the retinal surface, commonly related to age-related vitreous liquefaction. New flashes or floaters require dilated retinal examination because an associated retinal tear or detachment may be present.

\medterm{Retinal Tear}-- A full-thickness break in the neurosensory retina, often caused by vitreoretinal traction. It can permit fluid to enter the subretinal space and progress to rhegmatogenous retinal detachment; symptomatic tears are generally treated urgently with laser photocoagulation or cryotherapy.

\medterm{Central Retinal Artery Occlusion (CRAO)}-- Acute obstruction of the central retinal artery causing sudden, usually painless, severe monocular visual loss. It is an ocular stroke and requires urgent systemic and vascular evaluation because it may signal a high risk of cerebral ischemia.

\medterm{Retinal Vein Occlusion}-- Obstruction of a retinal vein producing retinal hemorrhages, edema, and variable visual loss. Branch and central forms may cause macular edema, retinal ischemia, and neovascular complications.

\medterm{Diabetic Retinopathy}-- A retinal microvascular complication of diabetes characterized by capillary leakage, occlusion, ischemia, and pathologic neovascularization. Nonproliferative disease may progress to proliferative retinopathy, macular edema, vitreous hemorrhage, and tractional retinal detachment.

\medterm{Central Serous Chorioretinopathy}-- Accumulation of serous fluid beneath the neurosensory retina, usually in the macular region, due to altered choroidal vascular permeability and retinal pigment-epithelium dysfunction. It often causes acute or subacute central visual distortion and may be associated with corticosteroid exposure.

\medterm{Keratoplasty}-- Surgical transplantation or reshaping of corneal tissue. Penetrating keratoplasty replaces the full corneal thickness, whereas lamellar techniques replace selected layers and may reduce rejection and recovery time.

\medterm{Corneal Edema}-- Excess fluid within the corneal stroma or epithelium causing corneal thickening, reduced transparency, glare, and blurred vision. Causes include endothelial dysfunction, acute glaucoma, inflammation, trauma, and postoperative injury.

\medterm{Ocular Surface Disease}-- A group of disorders affecting the cornea, conjunctiva, eyelids, and tear film, including dry eye, blepharitis, meibomian-gland dysfunction, and exposure-related disease. Symptoms and signs reflect tear-film instability, inflammation, epithelial injury, or abnormal lid function.


\medterm{Ophthalmic Emergency}-- An eye or visual condition requiring urgent evaluation to prevent irreversible vision loss or severe ocular injury. Examples include acute angle closure, retinal detachment, central retinal artery occlusion, endophthalmitis, chemical injury, orbital cellulitis, and open-globe trauma.

\medterm{Chemical Eye Injury}-- Ocular damage caused by exposure to an acidic or alkaline substance. Immediate copious irrigation, removal of particulate material, measurement of ocular-surface pH, and urgent ophthalmic assessment are required; alkali injuries may penetrate deeply and continue to worsen after exposure.
\medterm{Autoimmune keratitis}-- Corneal inflammation caused by the body's own immune system attacking corneal tissue, often associated with systemic autoimmune conditions such as rheumatoid arthritis, lupus, or granulomatosis with polyangiitis. It presents with corneal opacity, pain, and vision loss and requires immunosuppressive therapy.
\medterm{Binocular visual loss}-- Reduction in vision affecting both eyes simultaneously, typically resulting from bilateral optic neuropathy, chiasmal lesions, or diffuse bilateral retinal disease rather than monocular ocular pathology. Evaluation requires neuroimaging and visual field testing to localise the lesion along the visual pathway.
\medterm{Coloboma}-- A congenital defect caused by failure of the optic fissure to close during embryonic development, resulting in a keyhole-shaped gap in any ocular structure including the iris, retina, choroid, or optic nerve. Iris colobomas are the most common and appear as a comma-shaped defect in the inferior iris.
\medterm{Dislocation of the lens}-- Partial or complete displacement of the crystalline lens from its normal position within the pupillary axis due to disruption or weakness of the zonular fibers. Causes include Marfan syndrome, homocystinuria, trauma, and prior intraocular surgery; it may cause acute angle-closure glaucoma or blurred vision.
\medterm{Herpes zoster ophthalmicus}-- Reactivation of varicella-zoster virus latency in the trigeminal ganglion, with viral spread along the ophthalmic division (V1) causing painful vesicular skin eruption on the forehead, eyelid, and nose, with potential ocular involvement including keratitis, uveitis, and secondary glaucoma.
\medterm{IgG4-related ophthalmic disease}-- A fibroinflammatory condition affecting orbital and adnexal structures, characterised by lymphoplasmacytic infiltration enriched in IgG4-positive plasma cells, storiform fibrosis, and obliterative phlebitis. It commonly presents as lacrimal gland enlargement, orbital mass, or extraocular muscle involvement and responds to glucocorticoids.
\medterm{Inferior hemifield loss}-- Loss of vision in the lower half of the visual field in one or both eyes, which localises to lesions of the superior retina or superior optic radiations passing through the parietal lobe. Causes include superior retinal detachment, branch retinal artery occlusion, and parietal lobe stroke.
\medterm{Intraocular lens dislocation}-- Postoperative displacement of an implanted artificial intraocular lens from its intended position within the capsular bag or ciliary sulcus into the vitreous cavity or anterior chamber. It is most commonly caused by progressive zonular dehiscence and may require surgical repositioning or exchange.
\medterm{Intraocular tumor}-- Any neoplasm arising within the eye, with uveal melanoma being the most common primary malignancy in adults and retinoblastoma the most common in children. Diagnosis requires dilated fundus examination, ultrasound, and MRI; treatment depends on tumour type, size, and location.
\medterm{Iris atrophy}-- Thinning or loss of iris tissue resulting from trauma, chronic inflammation, ischaemia, or neurodegenerative disease, leading to a translucent or patched iris appearance. It may be sectoral or diffuse and can be associated with decreased pupillary response, photophobia, or secondary glaucoma.
\medterm{Lipemia retinalis}-- A creamy, pinkish-white discolouration of the retinal vessels caused by severely elevated serum triglycerides, typically exceeding 2,000 mg/dL. It is a sign of extreme hypertriglyceridaemia and indicates a risk for acute pancreatitis, requiring urgent lipid-lowering therapy and dietary management.
\medterm{Narrow angle glaucoma}-- Anatomical narrowing of the anterior chamber angle due to relative pupillary block or a shallow anterior chamber, creating risk for angle closure and acute intraocular pressure elevation. Identified on gonioscopy, it may be managed with laser peripheral iridotomy to equalise pressure between the posterior and anterior chambers.
\medterm{Necrotizing anterior scleritis}-- A severe, sight-threatening form of scleral inflammation characterised by ischaemic necrosis of the sclera with visible underlying uvea through thinned, grey-white scleral patches. It is strongly associated with systemic vasculitis and autoimmune disease and requires aggressive systemic immunosuppression to prevent scleral perforation and vision loss.
\medterm{Persistent pupillary membranes}-- Remnants of the fetal tunica vasculosa lentis that fail to regress after birth, appearing as fine pigmented strands bridging across the pupil or attaching from the iris collarette to the lens capsule. They are usually benign and asymptomatic but may rarely obstruct the visual axis requiring surgical or laser treatment.
\medterm{Prominent iris collarette}-- A visible thickening or ridge of the iris at the junction of the pupillary and ciliary zones, representing the site where the embryonic pupillary membrane was attached. It is typically a normal anatomical variant but can be a clinical sign in conditions such as Axenfeld-Rieger syndrome or iridocorneal endothelial syndrome.
\medterm{Radial keratotomy}-- A refractive surgical procedure, now largely obsolete, in which radial incisions are made in the peripheral corneal stroma to flatten the central cornea and reduce myopia. It was superseded by laser refractive surgery (LASIK, PRK) due to complications including diurnal visual fluctuation, glare, and hyperopic shift.
\medterm{Separation of the iris}-- Disinsertion or dialysis of the iris root from its attachment at the ciliary body, typically caused by blunt ocular trauma. It appears as a peripheral iridodialysis visible on slit-lamp examination and may cause monocular double vision, glare, or secondary glaucoma if extensive.
\medterm{Silicone oil hyperoleon}-- Emulsification of silicone oil tamponade used in vitreoretinal surgery, where the oil breaks into tiny droplets that migrate into the anterior chamber or trabecular meshwork. It can cause elevated intraocular pressure, band keratopathy, and chronic inflammation requiring oil removal and anterior chamber washout.
\medterm{Symblepharon}-- Adhesion of the palpebral conjunctiva (inner eyelid) to the bulbar conjunctiva or cornea, most commonly resulting from chemical burns, Stevens-Johnson syndrome, ocular cicatricial pemphigoid, or repeated surgical intervention. It restricts eye movement and may require surgical lysis with amniotic membrane or mucous membrane grafting.
\medterm{Traumatic bleb}-- A filtering bleb formed after glaucoma filtration surgery (trabeculectomy) that develops complications from trauma, or a thin-walled, elevated bleb that leaks or is at risk of rupture. Traumatic bleb failure or overfiltering can lead to hypotony, choroidal detachment, and endophthalmitis risk.
\medterm{Uveoparotid fever}-- A manifestation of sarcoidosis characterized by the triad of bilateral parotid gland enlargement, anterior uveitis, and fever, sometimes associated with facial nerve palsy. It reflects granulomatous inflammation and is treated with systemic corticosteroids and potentially disease-modifying agents.
\medterm{X-linked retinoschisis}-- An inherited vitreoretinal dystrophy caused by mutations in the RS1 gene, characterized by splitting of the retinal nerve fibre layer producing a spoke-wheel pattern of cystic spaces at the macula and peripheral retinal schisis. It primarily affects young males and may cause vitreous hemorrhage, retinal detachment, and progressive vision loss.
\medterm{Tractional retinal detachment}-- Separation of the neurosensory retina caused by contraction of fibrovascular membranes, commonly associated with proliferative diabetic retinopathy.

\medterm{Anophthalmia}-- Congenital absence of the eye due to failure of optic vesicle development, often associated with syndromes (e.g., SOX2 mutations). Diagnosis by absence of globe on imaging or clinical examination; management requires early socket expansion and prosthetic fitting for orbital and facial development.

\medterm{Microphthalmia}-- A congenitally small eye, ranging from mild (reduced axial length) to severe with structural anomalies. Associated with coloboma, congenital cataract, and genetic syndromes; can cause vision loss and requires optical correction, amblyopia therapy, and sometimes socket rehabilitation.

\medterm{Bietti's Crystalline Dystrophy}-- A rare autosomal recessive chorioretinal dystrophy characterized by tiny crystalline deposits in the retina and cornea with progressive chorioretinal atrophy, presenting with night blindness and progressive visual field loss beginning in early adulthood.

\medterm{Blepharospasm}-- Involuntary, bilateral tonic-clonic spasms of the eyelid muscles causing forced eye closure that can be functionally disabling. Most often essential (idiopathic) but can follow ocular irritation or brainstem lesions; first-line treatment is local botulinum toxin injection.

\medterm{Cerebral Visual Impairment (CVI)}-- Visual impairment caused by damage to the brain's visual pathways or occipital cortex rather than the eye itself; the leading cause of childhood visual impairment in developed countries. Features include poor visual attention, difficulty recognizing faces or objects, and variable vision; management is educational and rehabilitative.

\medterm{Convergence Insufficiency}-- Inability to maintain binocular convergence on near objects, causing eyestrain, headaches, blurred near vision, and transient diplopia during reading. Common in children and young adults; treated with convergence exercises and base-in prism, with surgery rarely needed.

\medterm{Low Vision}-- Visual impairment that cannot be fully corrected with spectacles, contact lenses, medication, or surgery, impairing daily activities. Managed with low-vision aids (magnifiers, telescopes, adaptive devices), lighting, and occupational therapy.

\medterm{Macular Hole}-- A full-thickness defect of the fovea usually caused by vitreomacular traction, presenting with central scotoma and distorted or reduced central vision. Staged by OCT; surgical vitrectomy with internal limiting membrane peel restores closure in most cases.

\medterm{Macular Pucker}-- An epiretinal membrane of fibrocellular tissue on the inner retinal surface that contracts to distort the macula, causing metamorphopsia and reduced vision. Diagnosed by OCT; vitrectomy with membrane peeling is reserved for visually significant symptoms.

\medterm{Ocular Histoplasmosis Syndrome (OHS)}-- Ocular disease following prior subclinical Histoplasma capsulatum infection, characterized by punched-out peripheral chorioretinal scars, peripapillary atrophy, and choroidal neovascularization that causes central vision loss. Active choroidal neovascularization is treated with anti-VEGF therapy.

\medterm{Presbyopia}-- Age-related loss of accommodation caused by lens hardening and loss of elasticity, beginning around age 40 and causing progressive difficulty with near vision. Corrected with reading glasses, bifocals, or progressive lenses; the most common refractive change of aging.

\medterm{Stargardt Disease}-- The most common inherited macular dystrophy, usually autosomal recessive and caused by ABCA4 gene mutations. Presents in childhood or adolescence with progressive central vision loss, retinal flecks, and bull's-eye maculopathy; no curative treatment exists.

\medterm{Usher Syndrome}-- An inherited disorder combining sensorineural hearing loss with retinitis pigmentosa, the most common cause of combined deaf-blindness. Three types are distinguished by severity and age of onset; management requires coordinated audiological, ophthalmological, and genetic care.

\medterm{Diplopia}-- See \hyperlink{Diplopia (Double Vision)}{Diplopia (Double Vision)}.

\medterm{Ectropion}-- Outward turning of the eyelid margin, most often the lower lid, exposing the conjunctiva and causing tearing, irritation, dryness, and incomplete eyelid closure. It may be involutional, paralytic, cicatricial, or mechanical; persistent exposure can damage the cornea.

\medterm{Fuchs Dystrophy}-- A progressive corneal endothelial disorder in which endothelial cell loss impairs fluid regulation and causes corneal oedema. Early disease may cause morning blur and glare; advanced disease produces persistent oedema, bullae, pain, and reduced vision. Diagnosis is clinical and supported by slit-lamp examination and corneal imaging.

\medterm{Iridocyclitis}-- Inflammation involving the iris and ciliary body, a form of anterior uveitis. It may cause eye pain, photophobia, redness, blurred vision, and cells or flare in the anterior chamber and may be associated with autoimmune disease, infection, trauma, or other causes.

\medterm{Nyctalopia}-- Difficulty seeing in dim light or darkness, commonly called night blindness. Causes include retinal disorders, vitamin-A deficiency, high refractive error, cataract, glaucoma, or medication and systemic disease affecting visual adaptation.

\medterm{Retinoschisis}-- Splitting of the retinal layers, producing a separation between inner and outer retinal tissue. It may be degenerative and peripheral or inherited and juvenile, and can cause visual-field loss, reduced acuity, or complications resembling retinal detachment.
